% Options for packages loaded elsewhere
\PassOptionsToPackage{unicode}{hyperref}
\PassOptionsToPackage{hyphens}{url}
\documentclass[
]{IEEEtran}
\usepackage{xcolor}
\usepackage{amsmath,amssymb}
\setcounter{secnumdepth}{-\maxdimen} % remove section numbering
\usepackage{iftex}
\ifPDFTeX
  \usepackage[T1]{fontenc}
  \usepackage[utf8]{inputenc}
  \usepackage{textcomp} % provide euro and other symbols
\else % if luatex or xetex
  \usepackage{unicode-math} % this also loads fontspec
  \defaultfontfeatures{Scale=MatchLowercase}
  \defaultfontfeatures[\rmfamily]{Ligatures=TeX,Scale=1}
\fi
\usepackage{lmodern}
\ifPDFTeX\else
  % xetex/luatex font selection
\fi
% Use upquote if available, for straight quotes in verbatim environments
\IfFileExists{upquote.sty}{\usepackage{upquote}}{}
\IfFileExists{microtype.sty}{% use microtype if available
  \usepackage[]{microtype}
  \UseMicrotypeSet[protrusion]{basicmath} % disable protrusion for tt fonts
}{}
\makeatletter
\@ifundefined{KOMAClassName}{% if non-KOMA class
  \IfFileExists{parskip.sty}{%
    \usepackage{parskip}
  }{% else
    \setlength{\parindent}{0pt}
    \setlength{\parskip}{6pt plus 2pt minus 1pt}}
}{% if KOMA class
  \KOMAoptions{parskip=half}}
\makeatother
\usepackage{,booktabs,array}
\newcounter{none} % for unnumbered tables
\usepackage{calc} % for calculating minipage widths
% Correct order of tables after \paragraph or \paragraph
\usepackage{etoolbox}
\makeatletter
% \patchcmd\longtable{\par}{\if@noskipsec\mbox{}\fi\par}{}{}
\makeatother
% Allow footnotes in longtable head/foot
\IfFileExists{footnotehyper.sty}{\usepackage{footnotehyper}}{\usepackage{footnote}}
% \makesavenoteenv{longtable}
\setlength{\emergencystretch}{3em} % prevent overfull lines
\providecommand{\tightlist}{%
  \setlength{\itemsep}{0pt}\setlength{\parskip}{0pt}}
\usepackage{bookmark}
\IfFileExists{xurl.sty}{\usepackage{xurl}}{} % add URL line breaks if available
\urlstyle{same}
\hypersetup{
  hidelinks,
  pdfcreator={LaTeX via pandoc}}

\author{}
\date{}

\begin{document}

\section{GNU RADIO TABANLI GÜÇ ETKİ ALANLI NOMA SİSTEMİ İLE ARDIŞIK
GİRİŞİM İPTALİ (SIC)
UYGULAMASI}\label{gnu-radio-tabanli-guxfcuxe7-etki-alanli-noma-sistemi-ile-ardiux15fik-giriux15fim-iptali-sic-uygulamasi}

\subsection{Bitirme Projesi Raporu}\label{bitirme-projesi-raporu}

\textbf{Bölüm:} Elektrik-Elektronik Mühendisliği\\
\textbf{Akademik Yıl:} 2025--2026

\begin{center}\rule{0.5\linewidth}{0.5pt}\end{center}

\subsection{İÇİNDEKİLER}\label{iuxe7indekiler}

\begin{enumerate}
\def\labelenumi{\arabic{enumi}.}
\tightlist
\item
  \hyperref[1-uxf6zet-abstract]{ÖZET (ABSTRACT)}
\item
  \hyperref[2-giriux15f-ve-literatuxfcr-uxf6zeti]{GİRİŞ VE LİTERATÜR
  ÖZETİ}
\item
  \hyperref[3-teorik-altyapux131-ve-matematiksel-model]{TEORİK ALTYAPI
  VE MATEMATİKSEL MODEL}
\item
  \hyperref[4-gnu-radio-tabanlux131-sistem-tasarux131mux131-blok-analizi]{GNU
  RADIO TABANLI SİSTEM TASARIMI (BLOK ANALİZİ)}
\item
  \hyperref[5-simuxfclasyon-ve-performans-deux11ferlendirmesi]{SİMÜLASYON
  VE PERFORMANS DEĞERLENDİRMESİ}
\item
  \hyperref[6-sonuuxe7-ve-gelecek-uxe7alux131ux15fmalar]{SONUÇ VE
  GELECEK ÇALIŞMALAR}
\item
  \hyperref[7-kaynakuxe7a]{KAYNAKÇA}
\item
  \hyperref[8-ekler]{EKLER}
\end{enumerate}

\begin{center}\rule{0.5\linewidth}{0.5pt}\end{center}

\subsection{1. ÖZET (ABSTRACT)}\label{uxf6zet-abstract}

Beşinci nesil (5G) ve sonrasındaki kablosuz haberleşme sistemlerinde
spektral verimliliğin artırılması, artan kullanıcı sayısı ve veri
trafiği karşısında en kritik mühendislik hedeflerinden biri hâline
gelmiştir. Bu bağlamda, Ortogonal Olmayan Çoklu Erişim (Non-Orthogonal
Multiple Access --- NOMA) tekniği, aynı zaman-frekans kaynağının birden
fazla kullanıcı tarafından eş zamanlı paylaşılmasına olanak tanıyan bir
fiziksel katman çözümü olarak öne çıkmaktadır.

Bu bitirme projesinde, Güç Etki Alanlı NOMA (Power-Domain NOMA ---
PD-NOMA) mimarisi kullanılarak iki kullanıcılı bir inen hat (downlink)
senaryosu için uçtan uca bir alıcı-verici (transceiver) sistemi
tasarlanmış ve GNU Radio Companion (GRC) ortamında gerçeklenmiştir.
Verici tarafında \textbf{Süperpozisyon Kodlaması (Superposition Coding)}
prensibi ile iki kullanıcının Diferansiyel BPSK (DBPSK) ile modüle
edilmiş sinyalleri farklı güç katsayıları (\(\alpha_1 = 0.894\),
\(\alpha_2 = 0.447\); sırasıyla \%80 ve \%20 güç tahsisi) atanarak aynı
taşıyıcı üzerinde birleştirilmiştir. Alıcı tarafında ise \textbf{Ardışık
Girişim İptali (Successive Interference Cancellation --- SIC)}
algoritması uygulanarak, önce baskın (güçlü) kullanıcının sinyali
çözülmüş, ardından bu sinyal yeniden oluşturulup bileşik sinyalden
çıkarılmak suretiyle zayıf kullanıcının verisi elde edilmiştir.

Sistem, IEEE 802.11n standardına uygun LDPC (Low-Density Parity-Check)
kanal kodlaması (\(n = 1296\), \(k = 648\), \(R = 1/2\)), kök
yükseltilmiş kosinüs (Root Raised Cosine --- RRC) darbe şekillendirme,
MMSE tabanlı polifaz sembol senkronizasyonu ve Costas döngüsü tabanlı
taşıyıcı kurtarma blokları ile donatılmıştır. AWGN (Additive White
Gaussian Noise) kanal modeli altında, \(f_s = 200 \text{ kHz}\)
örnekleme frekansı ve \(\text{sps} = 4\) sembol başına örnek
konfigürasyonunda gerçekleştirilen simülasyonlarda, Kullanıcı 1 (yakın
kullanıcı, \%80 güç tahsisi) için başarılı dosya transferi doğrulanmış
olup, çapraz korelasyon tabanlı SIC hizalama algoritması ile Kullanıcı 2
(uzak kullanıcı, \%20 güç tahsisi) sinyalinin çözümlenmesi
gerçekleştirilmiştir. Elde edilen sonuçlar, yazılım tabanlı radyo
(Software Defined Radio --- SDR) platformu üzerinde NOMA-SIC mimarisinin
uygulanabilirliğini ortaya koymakta ve gelecekte USRP E310 donanımı ile
gerçek zamanlı kablosuz doğrulamaya yönelik bir temel oluşturmaktadır.

\textbf{Anahtar Kelimeler:} NOMA, Süperpozisyon Kodlaması, Ardışık
Girişim İptali (SIC), GNU Radio, LDPC, DBPSK, SDR, 5G, Spektral
Verimlilik, USRP.

\begin{center}\rule{0.5\linewidth}{0.5pt}\end{center}

\subsection{2. GİRİŞ VE LİTERATÜR
ÖZETİ}\label{giriux15f-ve-literatuxfcr-uxf6zeti}

\subsubsection{2.1. Motivasyon ve Problem
Tanımı}\label{motivasyon-ve-problem-tanux131mux131}

Kablosuz haberleşme sistemlerinin evrimi, birinci nesilden (1G) beşinci
nesile (5G) ve ötesine uzanan süreçte, giderek artan veri trafiği,
kullanıcı yoğunluğu ve düşük gecikme gereksinimlerini karşılama
zorunluluğuyla şekillenmiştir. Uluslararası Telekomünikasyon Birliği
(ITU) tarafından tanımlanan IMT-2020 vizyonuna göre, 5G sistemlerinin
hedefleri arasında 10 Gbps tepe veri hızı, 1 ms gecikme süresi ve
\(10^6\) cihaz/km² bağlantı yoğunluğu yer almaktadır {[}1{]}. Bu
hedefler doğrultusunda, sınırlı frekans spektrumunun verimli
kullanılması kritik bir mühendislik problemi olarak öne çıkmaktadır.

Geleneksel kablosuz haberleşme sistemlerinde kullanılan Ortogonal Çoklu
Erişim (Orthogonal Multiple Access --- OMA) teknikleri --- FDMA, TDMA,
CDMA ve OFDMA --- kullanıcıları frekans, zaman, kod veya alt-taşıyıcı
boyutlarında birbirinden ayırarak ortogonallik sağlamaktadır. Bu
ortogonallik ilkesi, kullanıcılar arası girişimi (inter-user
interference) ortadan kaldırma avantajını sunarken, aynı zamanda kaynak
tahsisinde esnekliği ve toplam spektral verimliliği sınırlandırmaktadır.
Özellikle heterojen kullanıcı dağılımlarında --- yani kanal koşullarının
büyük farklılıklar gösterdiği senaryolarda --- OMA yöntemlerinin Shannon
kapasitesine yaklaşma kabiliyeti zayıflamaktadır.

\subsubsection{2.2. Non-Ortogonal Çoklu Erişim (NOMA)
Kavramı}\label{non-ortogonal-uxe7oklu-eriux15fim-noma-kavramux131}

NOMA, ortogonallik kısıtlamasını kaldırarak birden fazla kullanıcının
\textbf{aynı zaman-frekans kaynağını} eş zamanlı olarak paylaşmasını
mümkün kılan bir fiziksel katman çoklu erişim tekniğidir. Literatürde
NOMA teknikleri iki ana kategoride sınıflandırılmaktadır:

\begin{enumerate}
\def\labelenumi{\arabic{enumi}.}
\item
  \textbf{Güç Etki Alanlı NOMA (Power-Domain NOMA --- PD-NOMA):}
  Kullanıcılara farklı güç seviyeleri atanarak sinyallerin süperpozisyon
  kodlaması yoluyla üst üste bindirilmesine dayanır. Alıcı tarafında SIC
  algoritması ile girişim giderimi yapılır {[}2{]}.
\item
  \textbf{Kod Etki Alanlı NOMA (Code-Domain NOMA):} SCMA (Sparse Code
  Multiple Access) ve MUSA (Multi-User Shared Access) gibi teknikleri
  kapsayan, kullanıcıların seyrek veya düşük korelasyonlu kodlarla
  ayrıştırıldığı yaklaşımlardır.
\end{enumerate}

Bu projede, implementasyon karmaşıklığı ve kavramsal netlik açısından
\textbf{PD-NOMA} mimarisi tercih edilmiştir.

\subsubsection{2.3. OMA ve NOMA
Karşılaştırması}\label{oma-ve-noma-karux15fux131laux15ftux131rmasux131}

Aşağıdaki tablo, OMA ve NOMA tekniklerinin temel karakteristiklerini
karşılaştırmaktadır:

% {\def\LTcaptype{none} % do not increment counter
\begin{tabular{@{}
  >{\raggedright\arraybackslash}p{(\linewidth - 4\tabcolsep) * \real{0.3333}}
  >{\raggedright\arraybackslash}p{(\linewidth - 4\tabcolsep) * \real{0.3333}}
  >{\raggedright\arraybackslash}p{(\linewidth - 4\tabcolsep) * \real{0.3333}}@{}}
\toprule
\begin{minipage}[b]{\linewidth}\raggedright
Özellik
\end{minipage} & \begin{minipage}[b]{\linewidth}\raggedright
OMA (OFDMA vb.)
\end{minipage} & \begin{minipage}[b]{\linewidth}\raggedright
PD-NOMA
\end{minipage} \\
\midrule
\bottomrule
\textbf{Kaynak Tahsisi} & Her kullanıcıya ayrık alt-taşıyıcı/zaman
dilimi & Tüm kullanıcılar aynı kaynağı paylaşır \\
\textbf{Kullanıcılar Arası Girişim} & Ortogonallik ile sıfırlanır &
Kontrollü girişim, alıcıda SIC ile giderilir \\
\textbf{Spektral Verimlilik} & Kullanıcı sayısı ile ölçeklenmez &
Kullanıcı sayısı arttıkça kapasite kazancı sunar \\
\textbf{Kullanıcı Adaleti (Fairness)} & Zayıf kullanıcılar aleyhine
eşitsizlik & Güç tahsisi ile adalet kontrol edilebilir \\
\textbf{Alıcı Karmaşıklığı} & Düşük (doğrudan çözümleme) & Yüksek (SIC
algoritması gereklidir) \\
\textbf{Shannon Kapasitesine Yakınlık} & Sınırlı & Broadcast kanal
kapasitesine yaklaşır \\
\end{tabular
}

NOMA'nın en belirgin üstünlüğü, aynı spektral kaynak üzerinde birden
fazla kullanıcıya hizmet vermesi nedeniyle toplam sistem kapasitesini
artırmasıdır. Broadcast kanal kapasitesine ulaşmak için, güçlü
kullanıcıya atanan düşük güç ve zayıf kullanıcıya atanan yüksek güç ile
Shannon sınırına yaklaşılabilmektedir. Ancak, bu kazanım alıcı
tarafındaki SIC karmaşıklığı pahasına elde edilmektedir.

\subsubsection{2.4. Literatürde NOMA
Çalışmaları}\label{literatuxfcrde-noma-uxe7alux131ux15fmalarux131}

NOMA kavramının teorik temelleri, çok kullanıcılı bilgi teorisi
çerçevesinde Cover (1972) tarafından ortaya konulan yayın kanalı
(broadcast channel) kapasitesi analizine dayanmaktadır {[}3{]}. Saito
vd. (2013) tarafından yayımlanan çalışmada, güç etki alanlı NOMA'nın
3GPP-LTE sistemlerine entegrasyonu ilk kez önerilmiştir {[}4{]}. Dai vd.
(2015), NOMA'nın 5G kablosuz ağlar için bir anahtar teknoloji olduğunu
ortaya koyan kapsamlı bir derleme çalışması sunmuştur {[}5{]}. Ding vd.
(2014), çeşitli kanal koşulları altında PD-NOMA'nın ergodik kapasitesini
ve kesinti olasılığını (outage probability) analiz etmiştir {[}6{]}.
İslam vd. (2017) ise güç tahsisi stratejileri ve SIC alıcı tasarımına
ilişkin detaylı bir çerçeve geliştirmiştir {[}7{]}.

LDPC kodlarının NOMA sistemlerine entegrasyonu açısından, Gallager
(1962) tarafından ortaya konulan temel LDPC çerçevesi {[}8{]} ve MacKay
ve Neal (1996) tarafından gösterilen Shannon sınırına yakın performans
{[}9{]}, bu projede kullanılan IEEE 802.11n standardı LDPC kodlarının
teorik altyapısını oluşturmaktadır.

GNU Radio platformu üzerinde NOMA implementasyonu açısından, SDR
(Software Defined Radio) tabanlı çalışmalar sınırlı sayıdadır ve bu
proje, uçtan uca bir PD-NOMA transceiver'ın GNU Radio ortamında IEEE
802.11n LDPC kodlama ve DBPSK modülasyonu ile birlikte gerçeklenmesi
bakımından özgün bir katkı sunmaktadır.

\subsubsection{2.5. Projenin Amacı ve
Kapsamı}\label{projenin-amacux131-ve-kapsamux131}

Bu projenin temel amacı, iki kullanıcılı inen hat PD-NOMA senaryosunda:

\begin{enumerate}
\def\labelenumi{\arabic{enumi}.}
\tightlist
\item
  Verici tarafında süperpozisyon kodlaması ile sinyal birleştirme,
\item
  IEEE 802.11n LDPC kanal kodlaması (\(R = 1/2\)) ile hata dayanıklılığı
  sağlama,
\item
  AWGN kanal modeli altında (ayarlanabilir gürültü, frekans kayması ve
  zamanlama sapması ile) iletim simülasyonu,
\item
  Alıcı tarafında çapraz korelasyon tabanlı SIC algoritması ile ardışık
  kullanıcı çözümleme
\end{enumerate}

adımlarını kapsayan uçtan uca bir haberleşme zincirini GNU Radio
Companion ortamında tasarlamak, gerçeklemek ve performansını
değerlendirmektir.

\begin{center}\rule{0.5\linewidth}{0.5pt}\end{center}

\subsection{3. TEORİK ALTYAPI VE MATEMATİKSEL
MODEL}\label{teorik-altyapi-ve-matematiksel-model}

\subsubsection{3.1. Süperpozisyon Kodlaması (Superposition
Coding)}\label{suxfcperpozisyon-kodlamasux131-superposition-coding}

Güç etki alanlı NOMA'nın temel mekanizması, birden fazla kullanıcının
sinyallerinin farklı güç katsayılarıyla ağırlıklandırılarak toplam bir
sinyal olarak iletilmesidir. İki kullanıcılı bir inen hat senaryosunda,
baz istasyonu (BS) tarafından iletilen süperpoze sinyal aşağıdaki
şekilde ifade edilmektedir:

\[s(t) = \sqrt{a_1 P} \cdot x_1(t) + \sqrt{a_2 P} \cdot x_2(t)\]

Burada: - \(P\): Toplam iletim gücü, - \(x_1(t)\) ve \(x_2(t)\):
Sırasıyla Kullanıcı 1 (yakın/güçlü kanal) ve Kullanıcı 2 (uzak/zayıf
kanal) için modüle edilmiş sembol dizileri, - \(a_1\) ve \(a_2\): Güç
tahsis katsayıları (\(a_1 + a_2 = 1\), \(a_1 > a_2\)).

NOMA ilkesi gereğince, bu projede güç tahsisi:

\[a_1 = 0.8, \quad a_2 = 0.2\]

olarak belirlenmiştir. Normalize edilmiş birim güç (\(P = 1\)) varsayımı
altında, genlik katsayıları:

\[\alpha_1 = \sqrt{a_1} = \sqrt{0.8} \approx 0.894 \approx \frac{2}{\sqrt{5}}\]

\[\alpha_2 = \sqrt{a_2} = \sqrt{0.2} \approx 0.447 \approx \frac{1}{\sqrt{5}}\]

olarak hesaplanmıştır. Güç normalizasyonu doğrulaması:

\[\alpha_1^2 + \alpha_2^2 = 0.8 + 0.2 = 1.0\]

Süperpoze sinyal, genlik cinsinden:

\[s(t) = 0.894 \cdot x_1(t) + 0.447 \cdot x_2(t)\]

şeklinde yazılabilir. Burada Kullanıcı 1, yakın kullanıcı (near user)
olarak daha yüksek güç alırken; Kullanıcı 2, uzak kullanıcı (far user)
olarak daha düşük güç almaktadır.

\paragraph{3.1.1. Güç Tahsis Stratejisi ve Konstelasyon
Analizi}\label{guxfcuxe7-tahsis-stratejisi-ve-konstelasyon-analizi}

BPSK modülasyonu altında (\(x_i \in \{-1, +1\}\)), süperpoze sinyalin
konstelasyon uzayı dört noktadan oluşmaktadır:

% {\def\LTcaptype{none} % do not increment counter
\begin{tabular{@{}cccc@{}}
\toprule
\(x_1\) & \(x_2\) & \(s = \alpha_1 x_1 + \alpha_2 x_2\) & Bölge \\
\midrule
\bottomrule
+1 & +1 & \(+0.894 + 0.447 = +1.341\) & Pozitif \\
+1 & −1 & \(+0.894 − 0.447 = +0.447\) & Pozitif \\
−1 & +1 & \(−0.894 + 0.447 = −0.447\) & Negatif \\
−1 & −1 & \(−0.894 − 0.447 = −1.341\) & Negatif \\
\end{tabular
}

Bu konstelasyon yapısında, sıfır karar eşiği (\(\text{threshold} = 0\))
kullanılarak Kullanıcı 1'in (\(x_1\)) sembol kararı güvenilir bir
şekilde verilebilmektedir; zira \(\alpha_1 > \alpha_2\) koşulu, tüm
olası \(s\) değerlerinin \(x_1\)'in işaretiyle aynı yönde kalmasını
garanti etmektedir. Bu özellik, SIC'nin ilk aşamasında güçlü
kullanıcının doğru çözülmesinin temelini oluşturmaktadır.

\subsubsection{3.2. AWGN Kanal Modeli}\label{awgn-kanal-modeli}

Verici çıkışından alıcı girişine kadar olan iletim kanalı, Toplamsal
Beyaz Gauss Gürültüsü (AWGN) modeli ile temsil edilmektedir:

\[r(t) = h \cdot s(t) + n(t)\]

Burada: - \(h\): Kanal katsayısı (bu projede \(h = 1.0\), düz
sönümleme), - \(n(t)\): Sıfır ortalamalı ve \(\sigma_n^2\) varyanslı
karmaşık Gauss gürültüsü.

GNU Radio'daki \texttt{Channel\ Model} bloğunda konfigüre edilen
parametreler:

% {\def\LTcaptype{none} % do not increment counter
\begin{tabular{@{}
  >{\raggedright\arraybackslash}p{(\linewidth - 4\tabcolsep) * \real{0.3333}}
  >{\raggedright\arraybackslash}p{(\linewidth - 4\tabcolsep) * \real{0.3333}}
  >{\raggedright\arraybackslash}p{(\linewidth - 4\tabcolsep) * \real{0.3333}}@{}}
\toprule
\begin{minipage}[b]{\linewidth}\raggedright
Parametre
\end{minipage} & \begin{minipage}[b]{\linewidth}\raggedright
Değer
\end{minipage} & \begin{minipage}[b]{\linewidth}\raggedright
Açıklama
\end{minipage} \\
\midrule
\bottomrule
Gürültü Gerilimi (\(\sigma_n\)) & \(0.1\) (ayarlanabilir, 0--2 aralığı)
& AWGN standart sapması \\
Frekans Kayması (\(\Delta f\)) & \(0.01\) (ayarlanabilir, \(\pm 0.25\)
aralığı) & Normalize frekans ofseti \\
Zamanlama Sapması (\(\epsilon\)) & \(1.0001\) (ayarlanabilir,
0.999--1.001 aralığı) & Örnekleme frekansı sapma oranı \\
Kanal Katsayıları & \([1.0]\) & Tek-dokunuşlu düz sönümleme \\
Etiket Geçirme & \texttt{True} & Akış etiketlerini korur \\
\end{tabular
}

Varsayılan gürültü gerilimi (\(\sigma_n = 0.1\)) ile birim güç iletim
varsayımı altında Sinyal-Gürültü Oranı (SNR):

\[\text{SNR} = \frac{P}{\sigma_n^2} = \frac{1}{(0.1)^2} = 100 \quad (\approx 20 \text{ dB})\]

Qt GUI kaydırıcıları (slider) aracılığıyla gürültü gerilimi, frekans
kayması ve zamanlama sapması gerçek zamanlı olarak ayarlanabilmekte; bu
sayede sistemin farklı kanal koşullarındaki davranışı interaktif olarak
incelenebilmektedir.

\subsubsection{3.3. Ardışık Girişim İptali (Successive Interference
Cancellation ---
SIC)}\label{ardux131ux15fux131k-giriux15fim-iptali-successive-interference-cancellation-sic}

SIC, NOMA alıcısının temel bileşenidir ve çok kullanıcılı girişimin
kademeli olarak giderilmesini sağlar. İki kullanıcılı senaryoda SIC
süreci aşağıdaki adımlardan oluşmaktadır:

\paragraph{Adım 1: Güçlü Kullanıcının Çözülmesi (User 1
Decoding)}\label{adux131m-1-guxfcuxe7luxfc-kullanux131cux131nux131n-uxe7uxf6zuxfclmesi-user-1-decoding}

Alınan sinyal \(r(t)\) üzerinde, baskın güce sahip olan Kullanıcı 1'in
sinyali doğrudan çözümlenir. Kullanıcı 2'nin sinyali bu aşamada girişim
olarak ele alınır:

\[\hat{x}_1 = \text{Dec}\left\{ r(t) \right\} = \text{Dec}\left\{ \alpha_1 x_1(t) + \alpha_2 x_2(t) + n(t) \right\}\]

Kullanıcı 1 için Sinyal-Girişim-Artı-Gürültü Oranı (SINR):

\[\text{SINR}_1 = \frac{a_1 P}{a_2 P + \sigma_n^2} = \frac{0.8}{0.2 + 0.01} = \frac{0.8}{0.21} \approx 3.81 \quad (\approx 5.81 \text{ dB})\]

\paragraph{Adım 2: Güçlü Sinyalin Yeniden Oluşturulması (Signal
Reconstruction)}\label{adux131m-2-guxfcuxe7luxfc-sinyalin-yeniden-oluux15fturulmasux131-signal-reconstruction}

Çözülen \(\hat{x}_1\) verisi, verici tarafındaki işlem zinciri ile aynı
adımlardan (LDPC kanal kodlama, BPSK sembol eşleme, güç ölçekleme)
geçirilerek yeniden oluşturulur:

\[\hat{s}_1(t) = \alpha_1 \cdot \text{Mod}\left\{ \text{Enc}\left\{ \hat{x}_1 \right\} \right\}\]

Bu projede, yeniden oluşturma süreci LDPC kodlama → BPSK sembol haritası
→ \(\times 0.894\) genlik ölçekleme adımlarını kapsamaktadır.
\texttt{chunks\_to\_symbols} bloğu, çözülmüş bitleri doğrudan
\(\{-1.0 + 0j, +1.0 + 0j\}\) karmaşık BPSK sembollerine eşlemektedir.

\paragraph{Adım 3: Girişim İptali (Interference
Subtraction)}\label{adux131m-3-giriux15fim-iptali-interference-subtraction}

Yeniden oluşturulan güçlü kullanıcı sinyali, SIC hizalama bloğu
tarafından zamanlama ve faz açısından hizalandıktan sonra alınan bileşik
sinyalden çıkarılır:

\[r_{\text{clean}}(t) = r(t) - \hat{s}_1(t - \hat{\tau}) \cdot \hat{A} \cdot e^{j\hat{\phi}}\]

Burada \(\hat{\tau}\) tahmin edilen zamanlama gecikmesi, \(\hat{A}\)
genlik düzeltme faktörü ve \(\hat{\phi}\) faz ofseti tahminidir. Eğer
kod çözme hatasız (\(\hat{x}_1 = x_1\)) ve hizalama mükemmel ise:

\[r_{\text{clean}}(t) = \alpha_2 x_2(t) + n(t)\]

\paragraph{Adım 4: Zayıf Kullanıcının Çözülmesi (User 2
Decoding)}\label{adux131m-4-zayux131f-kullanux131cux131nux131n-uxe7uxf6zuxfclmesi-user-2-decoding}

Temizlenmiş sinyal üzerinde Kullanıcı 2'nin verisi çözümlenir:

\[\hat{x}_2 = \text{Dec}\left\{ r_{\text{clean}}(t) \right\}\]

Kullanıcı 2 için SNR (ideal SIC sonrası):

\[\text{SNR}_2 = \frac{a_2 P}{\sigma_n^2} = \frac{0.2}{0.01} = 20 \quad (\approx 13.01 \text{ dB})\]

\paragraph{3.3.1. SIC Hata Yayılımı (Error
Propagation)}\label{sic-hata-yayux131lux131mux131-error-propagation}

SIC sürecinin başarısı, Adım 1'deki kod çözme doğruluğuna kritik biçimde
bağlıdır. Kullanıcı 1'in hatalı çözülmesi durumunda, artık girişim
(residual interference) ortaya çıkar:

\[r_{\text{clean}}(t) = \alpha_2 x_2(t) + n(t) + \alpha_1 \left[ x_1(t) - \hat{x}_1(t) \right]\]

\(\alpha_1 \left[ x_1(t) - \hat{x}_1(t) \right]\) terimi hata yayılımını
temsil etmekte olup, \(\alpha_1 / \alpha_2 = 2\) olması nedeniyle her
bir hatalı sembol, Kullanıcı 2'nin sinyalini \(2:1\) oranında bastırma
potansiyeline sahiptir.

Benzer şekilde, zamanlama hizalamasındaki 1 sembol kayma durumunda artık
girişim:

\[r_{\text{clean}}[n] = r[n] - \alpha_1 x_1[n-1] = \alpha_1(x_1[n] - x_1[n-1]) + \alpha_2 x_2[n] + n[n]\]

biçimini alır. \(\alpha_1 = 2\alpha_2\) olduğundan, bu diferansiyel
girişim terimi Kullanıcı 2 sinyalini tamamen bastırabilir. Bu nedenle,
güçlü kanal kodlaması (LDPC) ve hassas zamanlama hizalaması büyük önem
taşımaktadır.

\subsubsection{3.4. LDPC Kanal
Kodlaması}\label{ldpc-kanal-kodlamasux131}

Düşük Yoğunluklu Parite Kontrol (Low-Density Parity-Check --- LDPC)
kodları, Shannon sınırına yakın performans sunan modern ileri hata
düzeltme (Forward Error Correction --- FEC) kodlarıdır {[}8, 9{]}. Bu
projede kullanılan LDPC kodu, IEEE 802.11n (Wi-Fi) standardından alınmış
olup aşağıdaki parametrelere sahiptir:

% {\def\LTcaptype{none} % do not increment counter
\begin{tabular{@{}ll@{}}
\toprule
Parametre & Değer \\
\midrule
\bottomrule
\textbf{Kod sözcüğü uzunluğu (\(n\))} & 1296 bit \\
\textbf{Mesaj uzunluğu (\(k\))} & 648 bit \\
\textbf{Parite bitleri (\(n - k\))} & 648 bit \\
\textbf{Kod oranı (\(R = k/n\))} & \(648/1296 = 0.5\) (tam yarı oran) \\
\textbf{Parite kontrol matrisi} &
\texttt{n\_1296\_k\_0648\_ieee.alist} \\
\textbf{Alt-matris boyutu (\(Z\))} & 54 (yarı-döngüsel yapı) \\
\textbf{Kodlayıcı} & \texttt{fec.ldpc\_encoder} (parite matrisi
tabanlı) \\
\textbf{Kod çözücü} & \texttt{fec.ldpc\_decoder} (maks. iterasyon:
50) \\
\textbf{Delme deseni (Puncture Pattern)} &
\texttt{\textquotesingle{}11\textquotesingle{}} (delme yok) \\
\end{tabular
}

\paragraph{3.4.1. IEEE 802.11n LDPC Matrisi ve Yarı-Döngüsel
Yapı}\label{ieee-802.11n-ldpc-matrisi-ve-yarux131-duxf6nguxfcsel-yapux131}

Kullanılan LDPC kodu, IEEE 802.11n standardında tanımlanan yarı-döngüsel
(quasi-cyclic) bir yapıya sahiptir. Parite kontrol matrisi
\(\mathbf{H}\), \(Z = 54\) boyutlu döngüsel permütasyon
alt-matrislerinden oluşmaktadır:

\[\mathbf{H} = \begin{bmatrix} \mathbf{P}_{0,0} & \mathbf{P}_{0,1} & \cdots & \mathbf{P}_{0,23} \\ \mathbf{P}_{1,0} & \mathbf{P}_{1,1} & \cdots & \mathbf{P}_{1,23} \\ \vdots & \vdots & \ddots & \vdots \\ \mathbf{P}_{11,0} & \mathbf{P}_{11,1} & \cdots & \mathbf{P}_{11,23} \end{bmatrix}\]

Burada her \(\mathbf{P}_{i,j}\), \(54 \times 54\) boyutunda bir döngüsel
permütasyon matrisi veya sıfır matrisidir. Toplam boyutlar:
\(648 \times 1296\) (\(12 \times 24\) alt-matris ızgarası).

Matrisin sütun ağırlıkları değişkenlik göstermektedir: ilk 54 sütun
yüksek ağırlığa (11'e kadar) sahipken, geri kalan sütunlar 2--4 ağırlık
aralığındadır. Bu yapı, düzensiz (irregular) LDPC kodlarının üstün hata
düzeltme performansını sağlamaktadır.

Geçerli bir LDPC kod sözcüğü \(\mathbf{c}\), aşağıdaki koşulu sağlar:

\[\mathbf{H} \cdot \mathbf{c}^T = \mathbf{0} \pmod{2}\]

\paragraph{3.4.2. Alist Dosya Formatı}\label{alist-dosya-formatux131}

LDPC kodunun yapısını tanımlayan parite kontrol matrisi, MacKay alist
formatında saklanmaktadır (\texttt{n\_1296\_k\_0648\_ieee.alist}, 44,030
bayt). Bu formatta:

\begin{itemize}
\tightlist
\item
  İlk satır: \(n = 1296\) (sütun sayısı) ve \(m = 648\) (satır sayısı),
\item
  İkinci satır: Maksimum sütun ağırlığı ve maksimum satır ağırlığı,
\item
  Sonraki satırlar: Her sütun ve satırdaki sıfır olmayan elemanların
  pozisyonları.
\end{itemize}

\paragraph{3.4.3. Yük Boyutu ve LDPC
Eşleşmesi}\label{yuxfck-boyutu-ve-ldpc-eux15fleux15fmesi}

Bu projede her paket \(\text{payload\_size} = 77\) bayt (\(= 616\) bit)
kullanıcı verisine sahiptir. CRC-32 eklenmesi (4 bayt) sonrası \(81\)
bayt (\(= 648\) bit) elde edilmekte olup, bu değer LDPC kodunun mesaj
uzunluğu \(k = 648\) ile birebir örtüşmektedir. LDPC kodlama sonrası her
paket \(n = 1296\) bit uzunluğundaki kod sözcüğüne dönüşür; bu parametre
SIC hizalama bloğundaki \texttt{payload\_size\ =\ 1296} sembol değeriyle
doğrudan ilişkilidir.

\subsubsection{3.5. Diferansiyel BPSK (DBPSK)
Modülasyonu}\label{diferansiyel-bpsk-dbpsk-moduxfclasyonu}

Bu projede verici tarafında \textbf{Diferansiyel BPSK (DBPSK)}
modülasyonu kullanılmaktadır. Diferansiyel kodlamada, bilgi verinin
mutlak faz değerinde değil, ardışık semboller arasındaki faz farkında
taşınır:

\[d_k = b_k \oplus d_{k-1}\]

Burada \(b_k\) giriş biti, \(d_k\) diferansiyel kodlanmış bit ve
\(\oplus\) XOR işlemidir. BPSK sembol eşlemesi:

\[x_k = 2d_k - 1 = \begin{cases} +1 & d_k = 1 \\ -1 & d_k = 0 \end{cases}\]

Diferansiyel kodlamanın avantajı, alıcıda mutlak faz referansı
gerektirmeden çözümleme yapılabilmesidir; bu özellik, Costas döngüsünün
\(180°\) faz belirsizliğini ortadan kaldırmaktadır.

Alıcı tarafında diferansiyel çözme, özel gömülü Python bloğu
(\texttt{epy\_block\_0}) ile yumuşak karar (soft-decision) formatında
gerçekleştirilmektedir:

\[\hat{b}_k^{\text{soft}} = -(y_k \cdot y_{k-1})\]

Burada \(y_k\) yumuşak konstelasyon çıkışıdır. Negatif işaret, polarite
eşleştirmesi için uygulanmaktadır. Yumuşak karar bilgisinin korunması,
aşağı yöndeki LDPC kod çözücünün performansı açısından kritik öneme
sahiptir.

AWGN kanalında kodlanmamış DBPSK'nın bit hata olasılığı (BER):

\[P_b = \frac{1}{2} e^{-E_b/N_0}\]

Bu ifade, tutarlı (coherent) BPSK'ya kıyasla yaklaşık 1 dB'lik bir
performans kaybına karşılık gelmektedir; ancak diferansiyel kodlamanın
sağladığı faz belirsizliği çözümü ve pratik alıcı tasarımı kolaylığı bu
kaybı telafi etmektedir.

\subsubsection{3.6. Darbe Şekillendirme: Kök Yükseltilmiş Kosinüs (RRC)
Filtresi}\label{darbe-ux15fekillendirme-kuxf6k-yuxfckseltilmiux15f-kosinuxfcs-rrc-filtresi}

Sembolleri arası girişimi (Inter-Symbol Interference --- ISI) önlemek ve
bant genişliğini kontrol etmek amacıyla, Kök Yükseltilmiş Kosinüs darbe
şekillendirme filtresi kullanılmıştır. RRC filtresinin frekans tepkisi:

\[H_{\text{RRC}}(f) = \begin{cases} \sqrt{T_s} & |f| \leq \frac{1-\beta}{2T_s} \\ \sqrt{\frac{T_s}{2}\left[1 + \cos\left(\frac{\pi T_s}{\beta}\left(|f| - \frac{1-\beta}{2T_s}\right)\right)\right]} & \frac{1-\beta}{2T_s} < |f| \leq \frac{1+\beta}{2T_s} \\ 0 & |f| > \frac{1+\beta}{2T_s} \end{cases}\]

Burada \(\beta\) aşırı bant genişliği (excess bandwidth / roll-off)
faktörüdür. Bu projede \(\beta = 0.35\) olarak konfigüre edilmiştir.
Filtre uzunluğu \(11 \times \text{sps} = 44\) dokunuş (tap) olarak
belirlenmiştir. Verici tarafındaki \texttt{Constellation\ Modulator}
bloğu dahili RRC darbe şekillendirme uygularken, alıcı tarafında ayrı
bir \texttt{FFT\ RRC\ Filter} bloğu eşleştirilmiş filtreleme işlevini
üstlenmektedir.

Verici ve alıcıdaki eşleştirilmiş RRC filtrelerinin kaskadı, ISI'yi
sıfırlayan Yükseltilmiş Kosinüs (RC) darbe biçimini oluşturur:

\[H_{\text{RC}}(f) = H_{\text{RRC}}(f) \cdot H_{\text{RRC}}(f) = |H_{\text{RRC}}(f)|^2\]

\begin{center}\rule{0.5\linewidth}{0.5pt}\end{center}

\subsection{4. GNU RADIO TABANLI SİSTEM TASARIMI (BLOK
ANALİZİ)}\label{gnu-radio-tabanli-sistem-tasarimi-blok-analizi}

Bu bölümde, GNU Radio Companion (GRC) ortamında tasarlanan NOMA-SIC
dosya transferi sisteminin tüm blokları, parametreleri ve sinyal akış
yolları detaylı olarak incelenmektedir. Sistem, üç ana katmandan
oluşmaktadır: Verici (Transmitter), Kanal Modeli (Channel Model) ve
Alıcı (Receiver/SIC). Toplam blok envanteri 60'ın üzerindedir ve 16
sanal sinyal yolu (virtual sink/source) çifti ile karmaşık sinyal akışı
yönetilmektedir.

\subsubsection{4.0. Sistem Değişkenleri (Variable
Blocks)}\label{sistem-deux11fiux15fkenleri-variable-blocks}

Akış diyagramında tanımlanan 14 adet global sistem değişkeni aşağıdaki
tabloda özetlenmiştir:

% {\def\LTcaptype{none} % do not increment counter
\begin{tabular{@{}
  >{\raggedright\arraybackslash}p{(\linewidth - 4\tabcolsep) * \real{0.3333}}
  >{\raggedright\arraybackslash}p{(\linewidth - 4\tabcolsep) * \real{0.3333}}
  >{\raggedright\arraybackslash}p{(\linewidth - 4\tabcolsep) * \real{0.3333}}@{}}
\toprule
\begin{minipage}[b]{\linewidth}\raggedright
Değişken Adı
\end{minipage} & \begin{minipage}[b]{\linewidth}\raggedright
Değer
\end{minipage} & \begin{minipage}[b]{\linewidth}\raggedright
Açıklama
\end{minipage} \\
\midrule
\bottomrule
\texttt{samp\_rate} & \(200{,}000\) Hz (\(200\) kHz) & Örnekleme
frekansı \\
\texttt{sps} & \(4\) & Sembol başına örnek sayısı \\
\texttt{payload\_size} & \(77\) bayt & Paket yük boyutu \\
\texttt{preamble\_size} & \(250\) bayt & Preambül boyutu \\
\texttt{postamble\_size} & \(8\) bayt & Postambül boyutu \\
\texttt{noise} & \(0.1\) (Qt GUI aralığı: 0--2) & AWGN gürültü
gerilimi \\
\texttt{freq\_offset} & \(0.01\) (Qt GUI aralığı: \(\pm 0.25\)) &
Normalize frekans ofseti \\
\texttt{time\_offset} & \(1.0001\) (Qt GUI aralığı: 0.999--1.001) &
Zamanlama sapma oranı \\
\texttt{constel} & \texttt{digital.constellation\_bpsk()} & BPSK
konstelasyon nesnesi \\
\texttt{hdr} &
\texttt{digital.header\_format\_default(default\_access\_code,\ 0)} &
Başlık format nesnesi \\
\texttt{preamble\_syms} & \([0\text{xC0}, 0\text{xAF}] \times 4\) → BPSK
sembolleri & Korelasyon preambül dizisi (64 sembol) \\
\texttt{ldpc\_enc} & \texttt{variable\_ldpc\_encoder\_def} (IEEE alist)
& LDPC kodlayıcı tanımı \\
\texttt{ldpc\_dec} & \texttt{variable\_ldpc\_decoder\_def} (maks. iter:
50) & Kullanıcı 1 LDPC kod çözücü \\
\texttt{ldpc\_dec\_2} & \texttt{variable\_ldpc\_decoder\_def} (maks.
iter: 50) & Kullanıcı 2 LDPC kod çözücü \\
\end{tabular
}

Sembol hızı ve bant genişliği:

\[R_s = \frac{f_s}{\text{sps}} = \frac{200{,}000}{4} = 50{,}000 \text{ sembol/s}\]

\[B = R_s \cdot (1 + \beta) = 50{,}000 \times 1.35 = 67{,}500 \text{ Hz}\]

\subsubsection{4.A. Verici (Transmitter)
Katmanı}\label{a.-verici-transmitter-katmanux131}

Verici katmanı, iki paralel kodlama zincirinden ve bir süperpozisyon
birleştirme bloğundan oluşmaktadır. Her iki zincir yapısal olarak özdeş
olmakla birlikte, bağımsız blok örnekleri ve farklı güç katsayıları ile
konfigüre edilmiştir.

\paragraph{4.A.1. Kullanıcı 1 Verici Zinciri (Yakın Kullanıcı --- Near
User)}\label{a.1.-kullanux131cux131-1-verici-zinciri-yakux131n-kullanux131cux131-near-user}

Kullanıcı 1'in verici zinciri aşağıdaki blok sırasını takip etmektedir:

\textbf{File Source → Stream to Tagged Stream → Stream CRC32 → Repack
Bits (8→1) → Additive Scrambler → FEC Extended Encoder (LDPC) → Tagged
Stream Multiply Length (×2) → Protocol Formatter + Tagged Stream Mux
{[}Preambül, Başlık, LDPC Yük, Postambül{]} → Constellation Modulator
(DBPSK) → Tag Gate → Multiply Const (\(\alpha_1 = 0.894\))}

\paragraph{4.A.1.1. File Source (Dosya
Kaynağı)}\label{a.1.1.-file-source-dosya-kaynaux11fux131}

% {\def\LTcaptype{none} % do not increment counter
\begin{tabular{@{}ll@{}}
\toprule
Parametre & Değer \\
\midrule
\bottomrule
Dosya & \texttt{bpsk\_transmit.txt} \\
Tekrar (Repeat) & \texttt{True} \\
\end{tabular
}

Bu blok, iletilecek veriyi (metin dosyası) bayt akışı olarak sisteme
beslemektedir. Dosya, her biri \texttt{payload\_size\ =\ 77} bayt olan
paketler hâlinde \texttt{0–9} rakamlarının tekrarından oluşmaktadır.
\texttt{Repeat\ =\ True} ayarı, dosya sonuna ulaşıldığında başa
dönülerek sürekli iletim yapılmasını sağlar.

\paragraph{4.A.1.2. Stream to Tagged
Stream}\label{a.1.2.-stream-to-tagged-stream}

% {\def\LTcaptype{none} % do not increment counter
\begin{tabular{@{}ll@{}}
\toprule
Parametre & Değer \\
\midrule
\bottomrule
Paket Uzunluğu & \texttt{payload\_size} = 77 bayt \\
Etiket Anahtarı & \texttt{packet\_len} \\
\end{tabular
}

Sürekli bayt akışını, her biri 77 baytlık etiketlenmiş paketlere
bölmektedir.

\paragraph{4.A.1.3. Stream CRC32 (Döngüsel Artıklık
Kontrolü)}\label{a.1.3.-stream-crc32-duxf6nguxfcsel-artux131klux131k-kontroluxfc}

% {\def\LTcaptype{none} % do not increment counter
\begin{tabular{@{}ll@{}}
\toprule
Parametre & Değer \\
\midrule
\bottomrule
Etiket Anahtarı & \texttt{packet\_len} \\
Kontrol Modu & \texttt{False} (ekleme modu) \\
\end{tabular
}

Her pakete 4 baytlık CRC-32 sağlama toplamı ekler. CRC eklenmesi sonrası
paket boyutu \(77 + 4 = 81\) bayta (\(= 648\) bit \(= k\)) yükselir. Bu
değer, LDPC kodunun mesaj uzunluğu ile birebir eşleşmektedir.

\paragraph{4.A.1.4. Repack Bits (Bit Yeniden Paketleme ---
8→1)}\label{a.1.4.-repack-bits-bit-yeniden-paketleme-81}

% {\def\LTcaptype{none} % do not increment counter
\begin{tabular{@{}ll@{}}
\toprule
Parametre & Değer \\
\midrule
\bottomrule
Giriş bit/bayt & \texttt{8} \\
Çıkış bit/bayt & \texttt{1} \\
Endianness & MSB \\
Etiketli & \texttt{True} \\
\end{tabular
}

Bayt düzeyindeki veriyi, karıştırıcı ve LDPC kodlayıcının beklediği tek
bit/bayt formatına dönüştürür.

\paragraph{4.A.1.5. Additive Scrambler (Toplamsal
Karıştırıcı)}\label{a.1.5.-additive-scrambler-toplamsal-karux131ux15ftux131rux131cux131}

% {\def\LTcaptype{none} % do not increment counter
\begin{tabular{@{}ll@{}}
\toprule
Parametre & Değer \\
\midrule
\bottomrule
Maske (Mask) & \texttt{0x8A} \\
Tohum (Seed) & \texttt{0x7F} \\
Uzunluk & \texttt{7} \\
Bayt başına bit & \texttt{1} \\
Sıfırlama Etiketi & \texttt{packet\_len} \\
\end{tabular
}

Doğrusal geri beslemeli kaydırma yazmacı (LFSR) tabanlı karıştırıcı, 7
bitlik bir LFSR kullanarak veri dizisindeki uzun tekrarları kırar.
\texttt{reset\_tag\_key\ =\ "packet\_len"} parametresi, her paket
başında LFSR'ın tohum değerine sıfırlanmasını sağlayarak, alıcıdaki
de-karıştırıcı ile senkronizasyonu garanti eder. Karıştırıcının
sağladığı faydalar:

\begin{itemize}
\tightlist
\item
  Zamanlama kurtarma döngüsü için yeterli geçiş yoğunluğu,
\item
  DC sapmasının önlenmesi,
\item
  Spektral yayılmanın düzleştirilmesi.
\end{itemize}

\paragraph{4.A.1.6. FEC Extended Encoder (LDPC
Kodlayıcı)}\label{a.1.6.-fec-extended-encoder-ldpc-kodlayux131cux131}

% {\def\LTcaptype{none} % do not increment counter
\begin{tabular{@{}ll@{}}
\toprule
Parametre & Değer \\
\midrule
\bottomrule
Kodlayıcı Nesnesi & \texttt{ldpc\_enc} \\
Parite Matrisi Dosyası & \texttt{n\_1296\_k\_0648\_ieee.alist} \\
Mesaj Uzunluğu (\(k\)) & \(648\) bit \\
Kod Sözcüğü Uzunluğu (\(n\)) & \(1296\) bit \\
Delme Deseni & \texttt{\textquotesingle{}11\textquotesingle{}} (delme
yok) \\
\end{tabular
}

LDPC kodlayıcı, 648 bitlik mesaj bloklarını 1296 bitlik kod sözcüklerine
dönüştürerek \(R = 0.5\) oranlı ileri hata düzeltme koruması ekler.
Kodlama işlemi sonrası paket uzunluğu tam iki katına çıkar.

\paragraph{4.A.1.7. Tagged Stream Multiply Length
(×2.0)}\label{a.1.7.-tagged-stream-multiply-length-2.0}

LDPC kodlamanın paket uzunluğunu iki katına çıkarması nedeniyle,
\texttt{packet\_len} etiket değerinin güncellenmesi gerekmektedir. Bu
blok, etiket değerini \(\times 2.0\) çarpanı ile doğru değere günceller.

\paragraph{4.A.1.8. Protocol Formatter (Protokol
Biçimlendirici)}\label{a.1.8.-protocol-formatter-protokol-biuxe7imlendirici}

% {\def\LTcaptype{none} % do not increment counter
\begin{tabular{@{}
  >{\raggedright\arraybackslash}p{(\linewidth - 2\tabcolsep) * \real{0.5000}}
  >{\raggedright\arraybackslash}p{(\linewidth - 2\tabcolsep) * \real{0.5000}}@{}}
\toprule
\begin{minipage}[b]{\linewidth}\raggedright
Parametre
\end{minipage} & \begin{minipage}[b]{\linewidth}\raggedright
Değer
\end{minipage} \\
\midrule
\bottomrule
Format Nesnesi & \texttt{hdr} =
\texttt{header\_format\_default(default\_access\_code,\ 0)} \\
\end{tabular
}

GNU Radio'nun standart erişim kodunu (64-bit) ve yük uzunluğu alanını
içeren bir başlık üretir. Erişim kodu, alıcıdaki
\texttt{Correlate\ Access\ Code} bloğu tarafından paket sınırlarının
tespitinde kullanılır.

\paragraph{4.A.1.9. Tagged Stream Mux (Etiketli Akış
Çoğullayıcı)}\label{a.1.9.-tagged-stream-mux-etiketli-akux131ux15f-uxe7oux11fullayux131cux131}

% {\def\LTcaptype{none} % do not increment counter
\begin{tabular{@{}lll@{}}
\toprule
Giriş & İçerik & Boyut \\
\midrule
\bottomrule
Port 0 & Preambül (\texttt{{[}0xC0,\ 0xAF{]}} tekrarlayan) & 250 bayt \\
Port 1 & Başlık (erişim kodu + yük uzunluğu) & Değişken \\
Port 2 & LDPC Kodlu Yük & 1296 bit (162 bayt) \\
Port 3 & Postambül (\texttt{{[}0xC0,\ 0xAF{]}} tekrarlayan) & 8 bayt \\
\end{tabular
}

Dört giriş portundan gelen verileri sıralı olarak birleştirerek çerçeve
(frame) yapısını oluşturur:

\[\text{Çerçeve} = [\underbrace{\text{Preambül}}_{250 \text{ B}}] \| [\underbrace{\text{Başlık}}_{\text{Değişken}}] \| [\underbrace{\text{LDPC Yük}}_{162 \text{ B}}] \| [\underbrace{\text{Postambül}}_{8 \text{ B}}]\]

Preambül ve postambül, \texttt{Vector\ Source} bloklarından
\texttt{{[}0xC0,\ 0xAF{]}} bayt dizisinin tekrarı olarak üretilmektedir.

\paragraph{4.A.1.10. Constellation Modulator (Konstelasyon Modülatörü
---
DBPSK)}\label{a.1.10.-constellation-modulator-konstelasyon-moduxfclatuxf6ruxfc-dbpsk}

% {\def\LTcaptype{none} % do not increment counter
\begin{tabular{@{}ll@{}}
\toprule
Parametre & Değer \\
\midrule
\bottomrule
Konstelasyon & BPSK (noktalar: \([-1-j, -1+j, 1+j, 1-j]\)) \\
Diferansiyel Kodlama & \texttt{True} \\
SPS & \texttt{4} \\
Aşırı Bant Genişliği (\(\beta\)) & \(0.35\) \\
\end{tabular
}

Bu bileşik blok, aşağıdaki işlemleri kapsüllemektedir: 1.
\textbf{Diferansiyel kodlama:} Bit bilgisini ardışık semboller
arasındaki faz farkına dönüştürür. 2. \textbf{Konstelasyon eşleme:}
Kodlanmış bitleri BPSK sembollerine eşler. 3. \textbf{Üst örnekleme
(Upsampling):} Her sembolün \(\text{sps} = 4\) kopyasını oluşturur. 4.
\textbf{RRC Darbe Şekillendirme:} \(\beta = 0.35\) roll-off faktörü ile
bant sınırlama ve ISI önleme.

\paragraph{4.A.1.11. Tag Gate (Etiket
Kapısı)}\label{a.1.11.-tag-gate-etiket-kapux131sux131}

Verici tarafındaki paket etiketlerinin kanal model bloğuna sızmasını
engelleyerek, alıcının yalnızca kendi senkronizasyon mekanizmasıyla
çalışmasını garanti eder.

\paragraph{4.A.1.12. Multiply Const (Sabit Çarpma --- Güç
Tahsisi)}\label{a.1.12.-multiply-const-sabit-uxe7arpma-guxfcuxe7-tahsisi}

% {\def\LTcaptype{none} % do not increment counter
\begin{tabular{@{}ll@{}}
\toprule
Parametre & Değer \\
\midrule
\bottomrule
Sabit (\(\alpha_1\)) & \(0.894\) \\
\end{tabular
}

Kullanıcı 1'e atanan güç katsayısı (\(a_1 = 0.8\)) doğrultusunda, modüle
edilmiş sinyalin genliği \(\alpha_1 = \sqrt{0.8} \approx 0.894\) ile
ölçeklenir.

\paragraph{4.A.2. Kullanıcı 2 Verici Zinciri (Uzak Kullanıcı --- Far
User)}\label{a.2.-kullanux131cux131-2-verici-zinciri-uzak-kullanux131cux131-far-user}

Kullanıcı 2'nin verici zinciri, Kullanıcı 1 ile \textbf{yapısal olarak
özdeştir}; yalnızca aşağıdaki parametrelerde farklılık göstermektedir:

% {\def\LTcaptype{none} % do not increment counter
\begin{tabular{@{}
  >{\raggedright\arraybackslash}p{(\linewidth - 2\tabcolsep) * \real{0.5000}}
  >{\raggedright\arraybackslash}p{(\linewidth - 2\tabcolsep) * \real{0.5000}}@{}}
\toprule
\begin{minipage}[b]{\linewidth}\raggedright
Blok
\end{minipage} & \begin{minipage}[b]{\linewidth}\raggedright
Farklılık
\end{minipage} \\
\midrule
\bottomrule
\textbf{File Source} & Kaynak dosya: \texttt{bpsk\_transmit\_2.txt} \\
\textbf{FEC Extended Encoder} & Kodlayıcı nesnesi:
\texttt{fec\_extended\_encoder\_0\_0} (aynı LDPC alist, ayrı örnek) \\
\textbf{Multiply Const} & Sabit (\(\alpha_2\)): \(0.447\)
(\(\sqrt{0.2}\)) \\
\end{tabular
}

\paragraph{4.A.3. Süperpozisyon Birleştirme (Add
Block)}\label{a.3.-suxfcperpozisyon-birleux15ftirme-add-block}

% {\def\LTcaptype{none} % do not increment counter
\begin{tabular{@{}ll@{}}
\toprule
Parametre & Değer \\
\midrule
\bottomrule
Tip & Karmaşık (complex) \\
Giriş Sayısı & \texttt{2} \\
Vektör Uzunluğu & \texttt{1} \\
\end{tabular
}

İki kullanıcının güç-ölçeklenmiş sinyalleri, \texttt{blocks\_add\_xx\_0}
bloğu ile toplanarak süperpoze sinyal oluşturulur:

\[s(t) = 0.894 \cdot x_1(t) + 0.447 \cdot x_2(t)\]

Süperpoze sinyal, \texttt{virtual\_sink\ "transmit"} üzerinden kanal
modeline yönlendirilir.

\subsubsection{4.B. Kanal Modeli (Channel
Model)}\label{b.-kanal-modeli-channel-model}

% {\def\LTcaptype{none} % do not increment counter
\begin{tabular{@{}
  >{\raggedright\arraybackslash}p{(\linewidth - 4\tabcolsep) * \real{0.3333}}
  >{\raggedright\arraybackslash}p{(\linewidth - 4\tabcolsep) * \real{0.3333}}
  >{\raggedright\arraybackslash}p{(\linewidth - 4\tabcolsep) * \real{0.3333}}@{}}
\toprule
\begin{minipage}[b]{\linewidth}\raggedright
Parametre
\end{minipage} & \begin{minipage}[b]{\linewidth}\raggedright
Değer
\end{minipage} & \begin{minipage}[b]{\linewidth}\raggedright
Açıklama
\end{minipage} \\
\midrule
\bottomrule
\textbf{Gürültü Gerilimi} & \texttt{noise} = \(0.1\) (varsayılan) & Qt
GUI kaydırıcısı ile ayarlanabilir (0--2) \\
\textbf{Frekans Kayması} & \texttt{freq\_offset} = \(0.01\) & Normalize
frekans ofseti (\(\pm 0.25\) aralığı) \\
\textbf{Zamanlama Sapması (\(\epsilon\))} & \texttt{time\_offset} =
\(1.0001\) & Örnekleme frekansı sapma oranı (0.999--1.001) \\
\textbf{Kanal Katsayıları} & \([1.0]\) & Tek-dokunuşlu düz sönümleme \\
\textbf{Etiket Geçirme} & \texttt{True} & Akış etiketlerini korur \\
\textbf{Gürültü Tohumu} & \(0\) & Tekrarlanabilir simülasyon \\
\end{tabular
}

Kanal modeli, gerçekçi bir iletim ortamını taklit etmek üzere süperpoze
sinyale üç tür bozulma eklemektedir:

\begin{enumerate}
\def\labelenumi{\arabic{enumi}.}
\tightlist
\item
  \textbf{AWGN Gürültüsü:} \(\sigma_n = 0.1\) (varsayılan) ile Gauss
  gürültüsü,
\item
  \textbf{Frekans Kayması:} \(\Delta f / f_s = 0.01\) normalize ofset,
  taşıyıcı frekansı sapmasını simüle eder,
\item
  \textbf{Zamanlama Sapması:} \(\epsilon = 1.0001\), verici ve alıcı
  osilatörleri arasındaki saat frekansı uyumsuzluğunu temsil eder.
\end{enumerate}

Qt GUI kaydırıcıları sayesinde bu parametreler gerçek zamanlı olarak
ayarlanabilmekte ve sistemin farklı bozulma koşullarındaki davranışı
interaktif olarak gözlemlenebilmektedir.

\subsubsection{4.C. Alıcı (Receiver) Katmanı ve SIC
Döngüsü}\label{c.-alux131cux131-receiver-katmanux131-ve-sic-duxf6nguxfcsuxfc}

Alıcı, dört ana alt-sisteme ayrılmaktadır: (1) Eşleştirilmiş Filtreleme
ve Senkronizasyon, (2) Kullanıcı 1 Çözümleme, (3) SIC Yeniden Oluşturma
ve Çıkarma, (4) Kullanıcı 2 Çözümleme.

\paragraph{4.C.1. Eşleştirilmiş Filtreleme ve Senkronizasyon
Blokları}\label{c.1.-eux15fleux15ftirilmiux15f-filtreleme-ve-senkronizasyon-bloklarux131}

\paragraph{FFT RRC Filter (Eşleştirilmiş
Filtre)}\label{fft-rrc-filter-eux15fleux15ftirilmiux15f-filtre}

% {\def\LTcaptype{none} % do not increment counter
\begin{tabular{@{}ll@{}}
\toprule
Parametre & Değer \\
\midrule
\bottomrule
Roll-off (\(\alpha\)) & \(0.35\) \\
Filtre Uzunluğu & \(11 \times \text{sps} = 44\) dokunuş \\
Örnekleme Frekansı & \(200{,}000\) Hz \\
Sembol Hızı & \(50{,}000\) sembol/s \\
Desimason & \(1\) (çıkış hızı değişmez) \\
\end{tabular
}

Alıcının ön ucunda yer alan bu blok, verici tarafındaki RRC darbe
şekillendirme filtresinin eşleştirilmiş karşılığıdır. FFT tabanlı
uygulama, uzun filtrelerde hesaplama verimliliği sağlar. Eşleştirilmiş
filtreleme sonrası toplam darbe tepkisi, Nyquist ISI kriterini sağlayan
Yükseltilmiş Kosinüs biçimini alır.

\paragraph{Symbol Synchronizer (Sembol Senkronizörü --- Polifaz Saat
Kurtarma)}\label{symbol-synchronizer-sembol-senkronizuxf6ruxfc-polifaz-saat-kurtarma}

% {\def\LTcaptype{none} % do not increment counter
\begin{tabular{@{}ll@{}}
\toprule
Parametre & Değer \\
\midrule
\bottomrule
TED Tipi & \texttt{SIGNAL\_TIMES\_SLOPE\_ML} (Sinyal × Eğim ML
tahmini) \\
SPS & \texttt{4} \\
Döngü Bant Genişliği & \(0.045\) \\
Sönümleme Faktörü & \(1.0\) \\
Maksimum Sapma & \(1.5\) \\
Çıkış SPS & \texttt{1} \\
İnterpolasyon Tipi & \texttt{IR\_MMSE\_8TAP} (MMSE-8 tap) \\
Polifaz Filtre Sayısı & \texttt{32} \\
TED Kazancı & \texttt{0.1} \\
\end{tabular
}

Polifaz filtre bankası tabanlı sembol senkronizörü, alınan sinyalin
optimum örnekleme anlarını belirleyerek sembol başına 4 örnekten (sps=4)
tek bir optimum örneğe (sps=1) desimasyonu gerçekleştirir.

\textbf{Zamanlama Hata Detektörü (TED):}
\texttt{SIGNAL\_TIMES\_SLOPE\_ML} yöntemi, Maximum Likelihood tabanlı
bir zamanlama tahmin algoritması olup, sinyalin anlık değeri ile
eğiminin çarpımından zamanlama hatası bilgisini türetir. Bu yöntem,
özellikle düşük SNR'da \texttt{Mueller\ and\ Muller} veya
\texttt{Zero-Crossing} algoritmalarına kıyasla daha gürbüz (robust)
performans sunmaktadır.

\textbf{MMSE İnterpolasyon:} 8-dokunuşlu Minimum Ortalama Kare Hatası
interpolatörü, polifaz filtre bankasının 32 alt-filtresi arasında
sürekli zamanlı örnekleme anı tahmini yaparak optimum sembol değerini
üretir.

\textbf{Döngü Dinamikleri:} İkinci dereceden döngü filtresi,
\(\omega_n = 0.045\) bant genişliği ve \(\zeta = 1.0\) sönümleme ile
konfigüre edilmiştir. Maksimum sapma \(1.5\) değeri, zamanlama
kaymasının \(\pm 1.5\) sembol aralığında izlenmesine olanak tanır.

\paragraph{Costas Loop (Costas Döngüsü --- Taşıyıcı Frekansı/Faz
Kurtarma)}\label{costas-loop-costas-duxf6nguxfcsuxfc-taux15fux131yux131cux131-frekansux131faz-kurtarma}

% {\def\LTcaptype{none} % do not increment counter
\begin{tabular{@{}ll@{}}
\toprule
Parametre & Değer \\
\midrule
\bottomrule
Döngü Bant Genişliği & \(2\pi / 100 \approx 0.0628\) rad \\
Derece & Konstelasyon noktası sayısına göre otomatik \\
\end{tabular
}

Costas döngüsü, taşıyıcı frekansındaki artık faz ve frekans sapmalarını
(\(\Delta f = 0.01\)) telafi eder. Döngü bant genişliği
(\(\omega_n = 2\pi/100\)), gürültü bastırma ve izleme hızı arasındaki
dengeyi optimize eder.

\paragraph{Correlation Estimator (Korelasyon
Tahmincisi)}\label{correlation-estimator-korelasyon-tahmincisi}

% {\def\LTcaptype{none} % do not increment counter
\begin{tabular{@{}ll@{}}
\toprule
Parametre & Değer \\
\midrule
\bottomrule
Preambül Sembolleri & \texttt{preamble\_syms} (64 BPSK sembol) \\
Eşik & \(0.7\) \\
Yöntem & \texttt{THRESHOLD\_ABSOLUTE} \\
SPS & \texttt{1} (sembol senkronizörü çıkışında) \\
Mark Delay & \texttt{len(preamble\_syms)\ -\ 1\ =\ 63} \\
\end{tabular
}

Bu blok, Costas Loop çıkışında sürgülü çapraz korelasyon yaparak
preambül dizisini (\([0\text{xC0}, 0\text{xAF}] \times 4\) → 64 BPSK
sembol) tespit eder. Korelasyon tepesi eşik değerini (\(0.7\)) aştığında
\texttt{time\_est} ve \texttt{corr\_est} etiketleri eklenir. Bu
etiketler, SIC hizalama bloğu tarafından paket sınırlarının
belirlenmesinde kullanılmaktadır.

\paragraph{4.C.2. Kullanıcı 1 Çözümleme Zinciri (SIC Aşama
1)}\label{c.2.-kullanux131cux131-1-uxe7uxf6zuxfcmleme-zinciri-sic-aux15fama-1}

Senkronizasyon bloklarından sonra, Kullanıcı 1'in verisi aşağıdaki
zincir ile çözümlenir:

\textbf{Costas Loop → \texttt{virtual\_sink\ "recovery"} → Constellation
Soft Decoder → EPY Block 0 (Yumuşak Diferansiyel Çözücü) → Correlate
Access Code → FEC Extended Decoder (LDPC) → Tagged Stream Multiply
Length (×0.5) → Additive Scrambler (de-karıştırma) → Repack Bits (1→8) →
Stream CRC32 (Kontrol) → File Sink}

\paragraph{Constellation Soft Decoder (Yumuşak Konstelasyon Kod
Çözücü)}\label{constellation-soft-decoder-yumuux15fak-konstelasyon-kod-uxe7uxf6zuxfccuxfc}

% {\def\LTcaptype{none} % do not increment counter
\begin{tabular{@{}ll@{}}
\toprule
Parametre & Değer \\
\midrule
\bottomrule
Konstelasyon & \texttt{constel} (BPSK) \\
Gürültü Gücü & \(-1\) (otomatik tahmin) \\
\end{tabular
}

Karmaşık sembolleri, BPSK konstelasyon haritasına göre yumuşak karar
(soft-decision) çıkışına dönüştürür. Sert karar (hard-decision) yerine
yumuşak karar kullanılması, LDPC kod çözücünün performansını önemli
ölçüde artırmaktadır; zira yumuşak bilgi, her bitin güvenilirlik
derecesini de taşır.

\paragraph{EPY Block 0 --- Yumuşak Diferansiyel Çözücü (Soft
Differential
Decoder)}\label{epy-block-0-yumuux15fak-diferansiyel-uxe7uxf6zuxfccuxfc-soft-differential-decoder}

% {\def\LTcaptype{none} % do not increment counter
\begin{tabular{@{}ll@{}}
\toprule
Parametre & Değer \\
\midrule
\bottomrule
Sınıf & \texttt{gr.sync\_block} \\
Giriş/Çıkış & \texttt{float32} / \texttt{float32} \\
Modülüs & \texttt{2} (BPSK modu) \\
\end{tabular
}

Bu özel gömülü Python bloğu (\texttt{NOMA\_epy\_block\_0.py}, 104
satır), diferansiyel kodlamanın tersini yumuşak karar formatında
gerçekleştirmektedir:

\[\hat{b}_k^{\text{soft}} = -(y_k \cdot y_{k-1})\]

Burada \(y_k\) yumuşak konstelasyon çıkışıdır. Çarpım işlemi,
diferansiyel fazı çözerken yumuşak karar büyüklüğünü korur. Negatif
işaret, \texttt{Correlate\ Access\ Code} bloğunun beklediği polarite
konvansiyonuna uyum sağlar (pozitif değer → bit `1' eşlemesi).

Blok, \texttt{self.prev} durum değişkeni ile çağrılar arası bellek
tutarak, ardışık \texttt{work()} çağrıları arasında diferansiyel
referansın sürdürülmesini sağlar.

\paragraph{Correlate Access Code --- Tag Stream (Erişim Kodu
Korelasyonu)}\label{correlate-access-code-tag-stream-eriux15fim-kodu-korelasyonu}

% {\def\LTcaptype{none} % do not increment counter
\begin{tabular{@{}ll@{}}
\toprule
Parametre & Değer \\
\midrule
\bottomrule
Erişim Kodu & GNU Radio varsayılan 64-bit erişim kodu \\
Eşik & \texttt{0} \\
Etiket Adı & \texttt{packet\_len} \\
\end{tabular
}

Yumuşak bit akışında erişim kodunu arar ve eşleşme bulunduğunda
\texttt{packet\_len} etiketi ekler.

\paragraph{FEC Extended Decoder (LDPC Kod
Çözücü)}\label{fec-extended-decoder-ldpc-kod-uxe7uxf6zuxfccuxfc}

% {\def\LTcaptype{none} % do not increment counter
\begin{tabular{@{}ll@{}}
\toprule
Parametre & Değer \\
\midrule
\bottomrule
Kod Çözücü & \texttt{ldpc\_dec} \\
Delme Deseni & \texttt{\textquotesingle{}11\textquotesingle{}} \\
Maksimum İterasyon & \texttt{50} \\
\end{tabular
}

LDPC kod çözücü, 1296 bitlik kod sözcüklerini 648 bitlik mesaj
bloklarına çözer. 50 iterasyon limiti, hesaplama maliyeti ile hata
düzeltme performansı arasındaki dengeyi gözetmektedir.

\paragraph{Tagged Stream Multiply Length
(×0.5)}\label{tagged-stream-multiply-length-0.5}

LDPC kod çözme sonrası paket uzunluğu yarıya düştüğünden
(\(1296 \to 648\) bit), \texttt{packet\_len} etiketinin güncellenmesi
gerekmektedir.

\paragraph{Stream CRC32 (Kontrol
Modu)}\label{stream-crc32-kontrol-modu}

% {\def\LTcaptype{none} % do not increment counter
\begin{tabular{@{}ll@{}}
\toprule
Parametre & Değer \\
\midrule
\bottomrule
Kontrol Modu & \texttt{True} \\
\end{tabular
}

CRC-32 sağlama toplamını doğrular; hatalı paketler düşürülür, doğru
paketler çıkışa iletilir. Hata durumunda \texttt{Message\ Debug} bloğuna
bildirim gönderilir.

\paragraph{File Sink (Dosya
Çıkışı)}\label{file-sink-dosya-uxe7ux131kux131ux15fux131}

% {\def\LTcaptype{none} % do not increment counter
\begin{tabular{@{}ll@{}}
\toprule
Parametre & Değer \\
\midrule
\bottomrule
Dosya & \texttt{bpsk\_receive.txt} \\
Ekleme Modu & \texttt{False} \\
\end{tabular
}

Çözümlenen Kullanıcı 1 verisi bu dosyaya yazılır.

\paragraph{4.C.3. SIC Mekanizması --- Yeniden Kodlama ve
Çıkarma}\label{c.3.-sic-mekanizmasux131-yeniden-kodlama-ve-uxe7ux131karma}

SIC'nin kritik aşaması, Kullanıcı 1'in başarıyla çözülen verisinin
yeniden oluşturularak alınan bileşik sinyalden çıkarılmasıdır. Bu süreç
iki alt-bileşenden oluşur: Yeniden Kodlama Zinciri ve SIC Hizalama
Bloğu.

\paragraph{SIC Yeniden Kodlama
Zinciri}\label{sic-yeniden-kodlama-zinciri}

\textbf{Çözülmüş User 1 Bitleri (\texttt{virtual\_source\ "user\_2"}) →
FEC Extended Encoder (LDPC SIC) → Chunks to Symbols (BPSK) → Multiply
Const (\(\alpha_1 = 0.894\)) → \texttt{virtual\_sink\ "user\_2\_sub"}}

% {\def\LTcaptype{none} % do not increment counter
\begin{tabular{@{}
  >{\raggedright\arraybackslash}p{(\linewidth - 4\tabcolsep) * \real{0.3333}}
  >{\raggedright\arraybackslash}p{(\linewidth - 4\tabcolsep) * \real{0.3333}}
  >{\raggedright\arraybackslash}p{(\linewidth - 4\tabcolsep) * \real{0.3333}}@{}}
\toprule
\begin{minipage}[b]{\linewidth}\raggedright
Blok
\end{minipage} & \begin{minipage}[b]{\linewidth}\raggedright
Parametre
\end{minipage} & \begin{minipage}[b]{\linewidth}\raggedright
Değer
\end{minipage} \\
\midrule
\bottomrule
\textbf{LDPC Encoder (SIC)} & Kodlayıcı &
\texttt{fec\_extended\_encoder\_0\_0\_0} (aynı IEEE alist) \\
\textbf{Chunks to Symbols} & Sembol Tablosu & \([-1.0 + 0j, +1.0 + 0j]\)
(BPSK) \\
\textbf{Multiply Const (SIC)} & Sabit & \(0.894\) (\(= \alpha_1\)) \\
\end{tabular
}

Yeniden kodlama zincirinde önemli bir tasarım kararı,
\texttt{Constellation\ Modulator} bloğu yerine doğrudan
\texttt{Chunks\ to\ Symbols} bloğunun kullanılmasıdır. Bu yaklaşım: -
Darbe şekillendirme ve üst örnekleme adımlarını atlayarak daha basit bir
yeniden oluşturma sağlar, - SIC hizalama bloğunun sembol düzeyinde
çalışmasını mümkün kılar.

\paragraph{EPY Block 1 --- Bağımsız NOMA SIC Hizalama Bloğu
(Decoupled NOMA SIC
Aligner)}\label{epy-block-1-baux11fux131msux131z-noma-sic-hizalama-bloux11fu-decoupled-noma-sic-aligner}

Bu proje kapsamında geliştirilen en kritik özel blok
\texttt{NOMA\_epy\_block\_1.py} dosyasında yer almaktadır (7.817 bayt,
170 satır). \texttt{gr.basic\_block} sınıfından türetilmiş bu blok,
\texttt{general\_work()} fonksiyonu ile çalışmaktadır.

\textbf{Portlar:}

% {\def\LTcaptype{none} % do not increment counter
\begin{tabular{@{}
  >{\raggedright\arraybackslash}p{(\linewidth - 6\tabcolsep) * \real{0.2500}}
  >{\raggedright\arraybackslash}p{(\linewidth - 6\tabcolsep) * \real{0.2500}}
  >{\raggedright\arraybackslash}p{(\linewidth - 6\tabcolsep) * \real{0.2500}}
  >{\raggedright\arraybackslash}p{(\linewidth - 6\tabcolsep) * \real{0.2500}}@{}}
\toprule
\begin{minipage}[b]{\linewidth}\raggedright
Port
\end{minipage} & \begin{minipage}[b]{\linewidth}\raggedright
Yön
\end{minipage} & \begin{minipage}[b]{\linewidth}\raggedright
Tip
\end{minipage} & \begin{minipage}[b]{\linewidth}\raggedright
Kaynak
\end{minipage} \\
\midrule
\bottomrule
\texttt{in\_rx} (Port 0) & Giriş & Karmaşık & Korelasyon Tahmincisi
çıkışı \\
\texttt{in\_tx1} (Port 1) & Giriş & Karmaşık & Yeniden oluşturulan User
1 sinyali \\
\texttt{out\_rx2} (Port 0) & Çıkış & Karmaşık & SIC sonrası artık User 2
sinyali \\
\end{tabular
}

\textbf{Konfigürasyon Parametreleri:}

% {\def\LTcaptype{none} % do not increment counter
\begin{tabular{@{}lll@{}}
\toprule
Parametre & Değer & Açıklama \\
\midrule
\bottomrule
\texttt{sample\_rate} & \(200{,}000\) & Örnekleme frekansı \\
\texttt{near\_user\_amplitude} & \(0.864\) & SIC genlik kalibrasyonu \\
\texttt{payload\_size} & \(1296\) & LDPC kod sözcüğü uzunluğu
(sembol) \\
\texttt{payload\_offset} & \(64\) & Başlık uzunluğu (sembol) \\
\texttt{search\_window} & \(32\) & Korelasyon arama penceresi \\
\end{tabular
}

\textbf{Algoritma --- Etiket-Bağımsız Çapraz Korelasyon SIC:}

\begin{enumerate}
\def\labelenumi{\arabic{enumi}.}
\item
  \textbf{Bağımsız Girdi Tamponlama (Decoupled Buffering):}
  \texttt{forecast()} fonksiyonu \texttt{{[}0,\ 0{]}} döndürerek GNU
  Radio zamanlayıcısına ``sıfır girdi ile çalışabilir'' mesajı verir. Bu
  kritik tasarım kararı, \texttt{in\_rx} ve \texttt{in\_tx1} portlarının
  bağımsız tüketilmesini sağlayarak zamanlayıcı kilitlenmesini
  (scheduler deadlock) önler. Gelen veriler anında dahili dairesel
  tamponlara (\texttt{buffer\_rx}, \texttt{buffer\_tx1}) aktarılır.
\item
  \textbf{Tampon Taşma Koruması:} \texttt{buffer\_rx} maksimum 50.000
  örnek, \texttt{buffer\_tx1} maksimum \(5 \times \text{payload\_size}\)
  örnek ile sınırlandırılmıştır.
\item
  \textbf{DBPSK Dönüşümü:} Yeniden oluşturulan BPSK sembolleri
  (\(\pm 1\)), DBPSK domenine dönüştürülür:
  \[\text{tx1\_diff}[k] = \prod_{i=0}^{k} (-\text{sgn}(\text{Re}(\text{tx1}[i])))\]
  Bu dönüşüm, verici tarafındaki diferansiyel kodlama ile eşleşme
  sağlar.
\item
  \textbf{Çapraz Korelasyon Hizalama:} \texttt{scipy.signal.correlate()}
  fonksiyonu ile RX tamponu ve DBPSK-dönüştürülmüş TX1 arasında
  korelasyon hesaplanır:
  \[C(\tau) = \sum_{n} r[n] \cdot \hat{s}_1^*[n - \tau]\] Korelasyon
  tepesinin konumu optimal zamanlama ofsetini verir.
\item
  \textbf{Genlik ve Faz Tahmini:} Korelasyon tepesinden:

  \begin{itemize}
  \tightlist
  \item
    Genlik: \(\hat{A} = |C_{\text{peak}}| / E_{\text{tx1}}\) (burada
    \(E_{\text{tx1}}\) TX1 enerjisi)
  \item
    Faz: \(\hat{\phi} = \angle C_{\text{peak}}\)
  \item
    Dinamik genlik ölçekleme: \(\hat{A} > 1.10\) ise \(1.5\times\)
    küçültme uygulanır
  \end{itemize}
\item
  \textbf{Girişim Çıkarma:}
  \[\text{buffer\_rx}[\text{aligned}] \mathrel{-}= \text{tx1\_diff} \cdot \hat{A} \cdot e^{j\hat{\phi}}\]
\item
  \textbf{Paket İzleme:} İlk paket için 1000--3000 örnek aralığında
  geniş arama; sonraki paketler için beklenen konumdan \(\pm 128\) örnek
  sapma ile dar arama. Paketler arası mesafe \(\approx 3408\) örnek.
\item
  \textbf{Güvenli Çıkış Akışı:} Yalnızca SIC işlemi tamamlanmış
  (çıkarılmış veya güvenli) örnekler çıkışa iletilir.
\item
  \textbf{Hata Ayıklama:} \texttt{debug\_sic.txt} dosyasına her paket
  için kayma, faz ve genlik bilgileri yazılır.
\end{enumerate}

\textbf{\texttt{basic\_block} vs.~\texttt{sync\_block} Tasarım Kararı:}

\texttt{sync\_block} yapısının neden kullanılmadığı, bu projedeki en
kritik mimari kararlardan biridir. \texttt{sync\_block}, her iki giriş
portundan eşit sayıda örnek hazır olmasını bekler. Ancak SIC
döngüsündeki geri besleme yapısında:

\begin{enumerate}
\def\labelenumi{\arabic{enumi}.}
\tightlist
\item
  \texttt{in\_rx} (Costas Loop çıkışı) gecikmesiz veri sağlar,
\item
  \texttt{in\_tx1} (yeniden oluşturulan sinyal) ise User 1'in çözülmesi,
  LDPC kod çözme ve yeniden kodlama gecikmesini barındırır (en az 1
  paketlik çerçeve gecikmesi),
\item
  \texttt{sync\_block}, \texttt{in\_tx1} verisi gelmeden \texttt{in\_rx}
  verisini tüketmez,
\item
  \texttt{in\_tx1} hattı ise \texttt{in\_rx} işlenmeden yeni veri
  üretemez.
\end{enumerate}

Bu dairesel bağımlılık, GNU Radio zamanlayıcısının \textbf{tamamen
kilitlenmesine} yol açar. \texttt{basic\_block} yapısı,
\texttt{consume(0,\ n)} ve \texttt{consume(1,\ m)} çağrılarıyla portları
bağımsız tüketerek bu sorunu çözmektedir.

\paragraph{4.C.4. Kullanıcı 2 Çözümleme Zinciri (SIC Aşama
2)}\label{c.4.-kullanux131cux131-2-uxe7uxf6zuxfcmleme-zinciri-sic-aux15fama-2}

Girişimi temizlenmiş sinyal üzerinde Kullanıcı 2'nin çözümlenmesi,
Kullanıcı 1 zinciri ile yapısal olarak benzer alıcı bloklarından
oluşmaktadır:

\textbf{EPY Block 1 Çıkışı → Constellation Soft Decoder (2.) → EPY Block
0\_0 (Yumuşak Diferansiyel Çözücü) → Correlate Access Code (2.) → FEC
Extended Decoder (LDPC, \texttt{ldpc\_dec\_2}) → Tagged Stream Multiply
Length (×0.5) → Additive Scrambler (de-karıştırma) → Repack Bits (1→8) →
Stream CRC32 (Kontrol) → File Sink (\texttt{bpsk\_receive\_2.txt})}

Tüm kod çözme blokları, Kullanıcı 1 zinciri ile aynı parametrelerde
konfigüre edilmiştir. Ayrı blok örnekleri (\texttt{\_0}, \texttt{\_0\_0}
son ekleri) kullanılarak iki zincirin bağımsız iç durum değişkenlerine
sahip olması sağlanmıştır:

% {\def\LTcaptype{none} % do not increment counter
\begin{tabular{@{}lll@{}}
\toprule
Blok & Kullanıcı 1 Örneği & Kullanıcı 2 Örneği \\
\midrule
\bottomrule
Soft Constellation Decoder & \texttt{\_cf\_0} & \texttt{\_cf\_0\_0} \\
Soft Diff Decoder (EPY) & \texttt{epy\_block\_0} &
\texttt{epy\_block\_0\_0} \\
Correlate Access Code & \texttt{\_ts\_0} & \texttt{\_ts\_0\_0} \\
LDPC Decoder & \texttt{ldpc\_dec} & \texttt{ldpc\_dec\_2} \\
File Sink & \texttt{bpsk\_receive.txt} &
\texttt{bpsk\_receive\_2.txt} \\
\end{tabular
}

\paragraph{4.C.5. Qt GUI Görselleştirme ve Kontrol
Blokları}\label{c.5.-qt-gui-guxf6rselleux15ftirme-ve-kontrol-bloklarux131}

Gerçek zamanlı sistem izleme ve parametre ayarı için aşağıdaki Qt
bileşenleri akış diyagramına eklenmiştir:

\textbf{Görselleştirme Blokları:}

% {\def\LTcaptype{none} % do not increment counter
\begin{tabular{@{}
  >{\raggedright\arraybackslash}p{(\linewidth - 4\tabcolsep) * \real{0.3333}}
  >{\raggedright\arraybackslash}p{(\linewidth - 4\tabcolsep) * \real{0.3333}}
  >{\raggedright\arraybackslash}p{(\linewidth - 4\tabcolsep) * \real{0.3333}}@{}}
\toprule
\begin{minipage}[b]{\linewidth}\raggedright
Blok
\end{minipage} & \begin{minipage}[b]{\linewidth}\raggedright
Bağlantı Noktası
\end{minipage} & \begin{minipage}[b]{\linewidth}\raggedright
Görüntülenen Bilgi
\end{minipage} \\
\midrule
\bottomrule
\textbf{QT GUI Time Sink} (``Received Samples'') & Kanal çıkışı & Alınan
bileşik sinyalin zaman gösterimi (1024 örnek) \\
\textbf{QT GUI Time Sink} (``Recovered Symbols'') & Costas Loop çıkışı &
Senkronize edilmiş User 1 sembolleri (256 örnek) \\
\textbf{QT GUI Time Sink} (``USER 2'') & SIC çıkışı & SIC sonrası User 2
sembolleri (256 örnek) \\
\end{tabular
}

\textbf{Kontrol Kaydırıcıları (Qt GUI Range):}

% {\def\LTcaptype{none} % do not increment counter
\begin{tabular{@{}llll@{}}
\toprule
Kaydırıcı & Değişken & Aralık & Adım \\
\midrule
\bottomrule
Gürültü & \texttt{noise} & \(0.0\) -- \(2.0\) & \(0.01\) \\
Frekans Kayması & \texttt{freq\_offset} & \(-0.25\) -- \(+0.25\) &
\(0.001\) \\
Zamanlama Sapması & \texttt{time\_offset} & \(0.999\) -- \(1.001\) &
\(0.0001\) \\
\end{tabular
}

\subsubsection{4.D. Sanal Sinyal Yolları (Virtual Sink/Source
Pairs)}\label{d.-sanal-sinyal-yollarux131-virtual-sinksource-pairs}

Akış diyagramının okunabilirliği ve yönetilebilirliği için 16 sanal
sinyal yolu çifti kullanılmıştır:

% {\def\LTcaptype{none} % do not increment counter
\begin{tabular{@{}ll@{}}
\toprule
Akış Kimliği & Açıklama \\
\midrule
\bottomrule
\texttt{tx\_packet\_1} & Kullanıcı 1 çerçevelenmiş TX paketi \\
\texttt{tx\_packet\_2} & Kullanıcı 2 çerçevelenmiş TX paketi \\
\texttt{header\_1} / \texttt{header\_2} & Protokol başlıkları \\
\texttt{ldpc\_payload\_1} / \texttt{ldpc\_payload\_2} & LDPC kodlu
yükler \\
\texttt{transmit} & Süperpoze NOMA sinyali \\
\texttt{channel} & Kanal çıkışı \\
\texttt{recovery} & Costas Loop çıkışı (User 1 doğrudan çözümleme) \\
\texttt{recovery\_2} & Korelasyon Tahmincisi çıkışı (SIC hizalama) \\
\texttt{rx\_payload} / \texttt{rx\_payload\_2} & Tespit edilmiş
yükler \\
\texttt{user\_2} & Çözülmüş User 1 bitleri (SIC yeniden kodlama
girişi) \\
\texttt{user\_2\_sub} & Yeniden oluşturulan User 1 sinyali (SIC
çıkarma) \\
\texttt{ldpc\_rx} / \texttt{ldpc\_rx\_2} & Son çözülmüş bitler \\
\end{tabular
}

\subsubsection{4.E. Çerçeve Yapısı (Frame
Structure)}\label{e.-uxe7eruxe7eve-yapux131sux131-frame-structure}

Her kullanıcı paketinin tam çerçeve yapısı aşağıdaki tabloda
özetlenmiştir:

% {\def\LTcaptype{none} % do not increment counter
\begin{tabular{@{}
  >{\raggedright\arraybackslash}p{(\linewidth - 6\tabcolsep) * \real{0.2500}}
  >{\raggedright\arraybackslash}p{(\linewidth - 6\tabcolsep) * \real{0.2500}}
  >{\raggedright\arraybackslash}p{(\linewidth - 6\tabcolsep) * \real{0.2500}}
  >{\raggedright\arraybackslash}p{(\linewidth - 6\tabcolsep) * \real{0.2500}}@{}}
\toprule
\begin{minipage}[b]{\linewidth}\raggedright
Segment
\end{minipage} & \begin{minipage}[b]{\linewidth}\raggedright
İçerik
\end{minipage} & \begin{minipage}[b]{\linewidth}\raggedright
Boyut
\end{minipage} & \begin{minipage}[b]{\linewidth}\raggedright
Sembol Sayısı
\end{minipage} \\
\midrule
\bottomrule
\textbf{Preambül} & \texttt{{[}0xC0,\ 0xAF{]}} tekrarlayan dizi & 250
bayt (2000 bit) & 256 sembol \\
\textbf{Başlık (Header)} & 64-bit erişim kodu + yük uzunluğu & Değişken
& \textasciitilde64 sembol \\
\textbf{Yük (Payload)} & LDPC kodlu kullanıcı verisi & 162 bayt (1296
bit) & 1296 sembol \\
\textbf{Postambül} & \texttt{{[}0xC0,\ 0xAF{]}} tekrarlayan dizi & 8
bayt (64 bit) & \textasciitilde64 sembol \\
\textbf{Toplam} & --- & --- & \textbf{\textasciitilde1680 sembol} \\
\end{tabular
}

Preambül deseni
\texttt{{[}0xC0,\ 0xAF{]}\ =\ {[}11000000,\ 10101111{]}}: - Zengin geçiş
yapısı sayesinde sembol senkronizörünün zamanlama kilitleme süresini
hızlandırır, - Costas döngüsünün faz referansını yakalamasını
kolaylaştırır, - Korelasyon tahmincisi için yüksek auto-korelasyon
tepesi sağlar.

Postambül ise çerçeve sonundaki RRC filtresinin kuyruk etkilerini
absorbe eder.

\begin{center}\rule{0.5\linewidth}{0.5pt}\end{center}

\subsection{5. SİMÜLASYON VE PERFORMANS
DEĞERLENDİRMESİ}\label{simuxfclasyon-ve-performans-deux11ferlendirmesi}

\subsubsection{5.1. Simülasyon Ortamı ve
Koşulları}\label{simuxfclasyon-ortamux131-ve-koux15fullarux131}

% {\def\LTcaptype{none} % do not increment counter
\begin{tabular{@{}
  >{\raggedright\arraybackslash}p{(\linewidth - 2\tabcolsep) * \real{0.5000}}
  >{\raggedright\arraybackslash}p{(\linewidth - 2\tabcolsep) * \real{0.5000}}@{}}
\toprule
\begin{minipage}[b]{\linewidth}\raggedright
Parametre
\end{minipage} & \begin{minipage}[b]{\linewidth}\raggedright
Değer
\end{minipage} \\
\midrule
\bottomrule
\textbf{Platform} & GNU Radio Companion 3.10.12+ \\
\textbf{İşletim Sistemi} & Linux (Arch Linux) \\
\textbf{Programlama Dili} & Python 3.8+ \\
\textbf{Simülasyon Türü} & Yazılımsal geri-döngü (software loopback) \\
\textbf{Kanal Modeli} & AWGN (\(\sigma_n = 0.1\), SNR ≈ 20 dB),
\(\Delta f = 0.01\), \(\epsilon = 1.0001\) \\
\textbf{Veri Kaynağı} & Metin dosyası (dosya transferi senaryosu) \\
\textbf{Veri İçeriği} & \texttt{0–9} rakamları, her biri
\texttt{payload\_size} kez tekrarlı \\
\textbf{Doğrulama Yöntemi} & Verici/alıcı dosya karşılaştırması
(\texttt{filecmp.cmp()}) \\
\textbf{Hedef Donanım} & USRP E310 (2.435 GHz ISM bandı, gelecek
çalışma) \\
\end{tabular
}

\subsubsection{5.2. Çerçeve Yapısının Senkronizasyon Kararlılığı
Üzerindeki
Etkisi}\label{uxe7eruxe7eve-yapux131sux131nux131n-senkronizasyon-kararlux131lux131ux11fux131-uxfczerindeki-etkisi}

Tasarlanan çerçeve yapısı (Preambül + Başlık + Yük + Postambül), alıcı
senkronizasyon zincirinin kararlılığı açısından kritik bir rol
üstlenmektedir:

\paragraph{5.2.1. Preambül ve Zamanlama
Kurtarma}\label{preambuxfcl-ve-zamanlama-kurtarma}

250 baytlık (\(= 2000\) bit \(\approx 256\) sembol) preambül segmenti,
\texttt{{[}0xC0,\ 0xAF{]}} tekrarlayan deseniyle Symbol Synchronizer
bloğunun zamanlama hata detektörü için zengin geçiş yapısına sahip bir
eğitim dizisi sağlamaktadır. \(\text{sps} = 4\) konfigürasyonunda
\(256 \times 4 = 1024\) örnek zamanlama edinim penceresi sunulmaktadır;
bu değer, pratikte güvenilir kilit süresini garanti etmektedir.

Korelasyon Tahmincisi bloğu, 64 BPSK sembolden oluşan
\texttt{preamble\_syms} dizisi ile preambül sonunu tespit eder.
\texttt{THRESHOLD\_ABSOLUTE} yöntemi ve \(0.7\) eşik değeri ile: -
Gürültülü ortamda kaçırılan paket olasılığı (missed detection) düşük
tutulmuş, - Yanlış alarm (false alarm) oranı kabul edilebilir seviyede
sınırlandırılmıştır.

\paragraph{5.2.2. Başlık ve Paket Sınırı
Algılama}\label{baux15flux131k-ve-paket-sux131nux131rux131-algux131lama}

GNU Radio'nun standart 64-bit erişim kodu, düşük oto-korelasyon yan
loblarına sahip bir dizi olarak seçilmiştir.
\texttt{Correlate\ Access\ Code\ –\ Tag\ Stream} bloğu ile paket sınırı
tespiti yapılmakta; yük uzunluğu alanı ise her paketin boyutunu
belirlemektedir.

\paragraph{5.2.3. Postambül ve Kuyruk
Etkileri}\label{postambuxfcl-ve-kuyruk-etkileri}

8 baytlık postambül segmenti, paket sonunda RRC darbe şekillendirme
filtresinin geçici tepkisinin (transient response) sönümlenmesi için
yeterli sembol süresi sağlamaktadır. Bu sayede ardışık paketler arasında
ISI oluşumu engellenmektedir.

\subsubsection{5.3. LDPC Kodlamanın SIC Hata Yayılımı Üzerindeki
Rolü}\label{ldpc-kodlamanux131n-sic-hata-yayux131lux131mux131-uxfczerindeki-roluxfc}

SIC mekanizmasının başarısı, büyük ölçüde birinci aşamadaki (Kullanıcı
1) kod çözme doğruluğuna bağlıdır. IEEE 802.11n LDPC kodlaması, bu
bağlamda birden fazla kritik işlev görmektedir:

\paragraph{5.3.1. Hata Düzeltme
Kapasitesi}\label{hata-duxfczeltme-kapasitesi}

\(R = 0.5\) oranlı LDPC kodu, her 648 bitlik bilgi bloğuna 648 bitlik
artıklık (redundancy) eklemektedir. Yarı-döngüsel yapısı ve düzensiz
sütun ağırlıkları sayesinde, bu kod Shannon sınırına 1 dB mesafede
performans sunmaktadır. 50 iterasyonluk kod çözme limiti ile:

\begin{itemize}
\tightlist
\item
  Kullanıcı 1'in neredeyse hatasız çözülmesi (\(\text{BER} < 10^{-6}\))
  sağlanmakta,
\item
  SIC çıkarma işlemindeki artık girişimin
  (\(\alpha_1[x_1 - \hat{x}_1]\)) minimize edilmesi mümkün
  kılınmaktadır.
\end{itemize}

\paragraph{5.3.2. Yumuşak Karar ve LDPC
Sinerjisi}\label{yumuux15fak-karar-ve-ldpc-sinerjisi}

Alıcı zincirinde \texttt{Constellation\ Soft\ Decoder} →
\texttt{Soft\ Diff\ Decoder} yolu ile korunan yumuşak karar (LLR)
bilgisi, LDPC kod çözücünün iteratif mesaj geçirme (message passing)
algoritmasının tam potansiyeline ulaşmasını sağlamaktadır. Sert karar
tabanlı bir alıcıya kıyasla, yumuşak karar LDPC kombinasyonu tipik
olarak 2--3 dB SNR kazancı sunmaktadır.

\paragraph{5.3.3. CRC ile Çift Katmanlı
Koruma}\label{crc-ile-uxe7ift-katmanlux131-koruma}

CRC-32 kontrolü, LDPC kod çözme sonrası paket bütünlüğünü doğrulayan
ikinci bir koruma katmanı oluşturmaktadır. Hatalı paketlerin SIC yeniden
kodlama zincirine girmesi engellenerek, hata yayılımının kaskat etkisi
(cascade effect) sınırlandırılmaktadır.

\subsubsection{5.4. Gözlemlenen Sonuçlar ve
Tartışma}\label{guxf6zlemlenen-sonuuxe7lar-ve-tartux131ux15fma}

\paragraph{5.4.1. Kullanıcı 1 (Yakın Kullanıcı)
Performansı}\label{kullanux131cux131-1-yakux131n-kullanux131cux131-performansux131}

\(\text{SNR} \approx 20\) dB koşullarında, Kullanıcı 1 verici dosyası
(\texttt{bpsk\_transmit.txt}) ile alıcı çıkış dosyası
(\texttt{bpsk\_receive.txt}) arasında başarılı dosya transferi
doğrulanmıştır. \texttt{filecmp.cmp()} fonksiyonu ile otomatik doğrulama
gerçekleştirilmiştir. LDPC kodlamasının sağladığı hata düzeltme
kapasitesi, AWGN gürültüsü ve frekans/zamanlama sapmalarının etkisini
büyük ölçüde telafi etmiştir.

\paragraph{5.4.2. Kullanıcı 2 (Uzak Kullanıcı) ve SIC
Performansı}\label{kullanux131cux131-2-uzak-kullanux131cux131-ve-sic-performansux131}

SIC mekanizması üzerinden Kullanıcı 2 sinyalinin çözümlenmesi
gerçekleştirilmiştir. Proje geliştirme sürecinde aşağıdaki teknik
zorluklar gözlemlenmiş ve çözüm stratejileri uygulanmıştır:

\begin{enumerate}
\def\labelenumi{\arabic{enumi}.}
\item
  \textbf{Zamanlama Hizalama Hassasiyeti:} SIC çıkarma işleminde
  orijinal alınan sinyal ile yeniden oluşturulan Kullanıcı 1 sinyali
  arasındaki zamanlama hizalaması, örnek düzeyinde (\(< 1\) sembol)
  hassasiyet gerektirmektedir. Çapraz korelasyon tabanlı dinamik
  hizalama algoritması bu soruna çözüm olarak geliştirilmiştir.
\item
  \textbf{Genlik ve Faz Kalibrasyonu:} Kanal etkilerinin (faz dönmesi,
  genlik değişimi) yeniden oluşturulan sinyale doğru yansıtılması
  gerekmektedir. SIC hizalama bloğu, korelasyon tepesinden genlik ve faz
  tahminleri çıkararak dinamik kalibrasyon uygulamaktadır.
\item
  \textbf{Scheduler Kilitlenmesi:} \texttt{sync\_block} tabanlı ilk
  tasarım, geri besleme döngüsünün dairesel bağımlılığı nedeniyle GNU
  Radio zamanlayıcısının kilitlenmesine yol açmıştır.
  \texttt{basic\_block} yapısına geçiş ile bu sorun çözülmüştür.
\item
  \textbf{Sinyal Sızıntısı (Signal Bleeding):} 1 sembol hizalama hatası
  durumunda, \(\alpha_1/\alpha_2 = 2\) oranındaki güç farkı nedeniyle
  User 1 kalıntısının User 2 sinyalini bastırması gözlemlenmiştir.
  \texttt{corr\_est} bloğunun \texttt{mark\_delay} parametresinin doğru
  kalibrasyonu ile bu etki minimize edilmiştir.
\end{enumerate}

\subsubsection{5.5. Sistem Doğrulama
Yöntemleri}\label{sistem-doux11frulama-yuxf6ntemleri}

Sistemin doğrulanmasında aşağıdaki yöntemler kullanılmıştır:

\begin{enumerate}
\def\labelenumi{\arabic{enumi}.}
\tightlist
\item
  \textbf{Otomatik Dosya Karşılaştırması:} \texttt{filecmp.cmp()}
  fonksiyonu ile verici ve alıcı metin dosyalarının bayt düzeyinde
  karşılaştırması.
\item
  \textbf{Gerçek Zamanlı Zaman Domeni İzleme:} Üç adet Qt GUI Time Sink
  ile:

  \begin{itemize}
  \tightlist
  \item
    Alınan bileşik sinyal (kanal çıkışı)
  \item
    Senkronize edilmiş User 1 sembolleri
  \item
    SIC sonrası User 2 sembolleri
  \end{itemize}
\item
  \textbf{Mesaj Hata Ayıklama:} CRC başarısız paketlerin
  \texttt{Message\ Debug} bloğu üzerinden izlenmesi.
\item
  \textbf{SIC Hata Ayıklama Günlüğü:} \texttt{debug\_sic.txt} dosyası
  üzerinden her paket için zamanlama kayması, faz ofseti ve genlik
  kalibrasyonu bilgilerinin kaydedilmesi.
\item
  \textbf{Çevrimdışı Sinyal Analizi:} \texttt{rx\_chunk.npy} ve
  \texttt{tx1\_chunk.npy} dosyaları üzerinden numpy tabanlı genlik ve
  korelasyon analizi (\texttt{capture\_sic\_amplitude.py},
  \texttt{offline\_analysis.py}, \texttt{test\_correlation.py}).
\end{enumerate}

\begin{center}\rule{0.5\linewidth}{0.5pt}\end{center}

\subsection{6. SONUÇ VE GELECEK
ÇALIŞMALAR}\label{sonuuxe7-ve-gelecek-uxe7aliux15fmalar}

\subsubsection{6.1. Elde Edilen Bulguların
Özeti}\label{elde-edilen-bulgularux131n-uxf6zeti}

Bu bitirme projesinde, Güç Etki Alanlı NOMA (PD-NOMA) mimarisinin GNU
Radio Companion platformu üzerinde uçtan uca bir gerçeklemesi başarıyla
tamamlanmıştır. Projenin temel katkıları aşağıdaki şekilde
özetlenebilir:

\begin{enumerate}
\def\labelenumi{\arabic{enumi}.}
\item
  \textbf{Süperpozisyon Kodlaması Gerçeklemesi:} İki kullanıcının DBPSK
  modüleli sinyallerinin \(\alpha_1 = 0.894\) (\%80 güç) ve
  \(\alpha_2 = 0.447\) (\%20 güç) katsayıları ile birleştirilmesi
  başarıyla uygulanmıştır.
\item
  \textbf{IEEE 802.11n LDPC Kodlu İletim Zinciri:} \(R = 0.5\) oranlı
  LDPC kodlaması (\(n = 1296\), \(k = 648\)), Toplamsal Karıştırıcı
  (LFSR \(0\text{x8A}/0\text{x7F}\)), CRC-32, diferansiyel kodlama ve
  \(\beta = 0.35\) RRC darbe şekillendirme ile entegre edilmiş tam bir
  verici-alıcı zinciri oluşturulmuştur.
\item
  \textbf{SIC Mekanizması:} Güçlü kullanıcının (User 1) yumuşak karar
  LDPC kod çözme ile çözülmesi, çözülen verisinin LDPC ile yeniden
  kodlanıp BPSK sembollerine eşlenmesi, ve çapraz korelasyon tabanlı
  hizalama ile bileşik sinyalden çıkarılması işlemleri gerçeklenmiştir.
  \texttt{basic\_block} tabanlı bağımsız tamponlama yapısı ile
  zamanlayıcı kilitlenmesi sorunu çözülmüştür.
\item
  \textbf{Dosya Transferi Doğrulaması:} AWGN kanalı
  (\(\text{SNR} \approx 20\) dB, \(\Delta f = 0.01\),
  \(\epsilon = 1.0001\)) altında Kullanıcı 1 için başarılı metin dosyası
  transferi otomatik olarak doğrulanmıştır.
\item
  \textbf{Gerçek Zamanlı İzleme ve Kontrol:} Qt GUI Time Sink blokları
  ile zaman domeni görselleştirme ve kaydırıcılarla interaktif parametre
  ayarlama arayüzü oluşturulmuştur.
\item
  \textbf{Evrimsel Geliştirme Süreci:} Proje, basit BPSK dosya
  transferinden (LDPC'siz) → tek kullanıcılı LDPC-BPSK → iki kullanıcılı
  NOMA-SIC mimarisine doğru sistematik bir geliştirme sürecini
  izlemiştir. Bu yaklaşım, her aşamada doğrulama yapılmasını ve
  sorunların kademeli olarak çözülmesini mümkün kılmıştır.
\end{enumerate}

\subsubsection{6.2. Projenin Akademik ve Mühendislik
Katkısı}\label{projenin-akademik-ve-muxfchendislik-katkux131sux131}

Bu çalışma, NOMA'nın teorik çerçevesinin ötesine geçerek, SDR platformu
üzerinde tüm sinyal işleme blokları --- diferansiyel modülasyon, IEEE
802.11n LDPC kanal kodlama, RRC darbe şekillendirme, MMSE tabanlı sembol
senkronizasyonu, Costas döngüsü taşıyıcı kurtarma ve çapraz korelasyon
tabanlı SIC --- ile birlikte çalışan uçtan uca bir demonstrasyon
sunmaktadır.

Özellikle SIC hizalama bloğunun (\texttt{epy\_block\_1}) tasarımı, GNU
Radio'nun akış tabanlı (streaming) mimarisi içinde geri beslemeli bir
yapının nasıl kurulacağına dair özgün bir mühendislik çözümü ortaya
koymaktadır. \texttt{basic\_block} vs.~\texttt{sync\_block} seçiminin
zamanlayıcı kilitlenmesi üzerindeki etkisi, DBPSK dönüşüm algoritması ve
korelasyon tabanlı dinamik genlik/faz kalibrasyonu, bu projede ele
alınan özgün teknik problemlerdir.

\subsubsection{6.3. Gelecek
Çalışmalar}\label{gelecek-uxe7alux131ux15fmalar}

Bu projenin sonuçları, aşağıdaki gelecek çalışma yönlerini açmaktadır:

\begin{enumerate}
\def\labelenumi{\arabic{enumi}.}
\item
  \textbf{USRP E310 Donanım Doğrulaması:} Tasarlanan sistemin USRP E310
  (2.435 GHz ISM bandı) üzerinde gerçek zamanlı kablosuz iletim ile
  doğrulanması. Proje kapsamında oluşturulan
  \texttt{Bpsk\_file\_transfer\_loopback.grc} akış diyagramı, bu
  aşamanın temelini oluşturmaktadır. Gerçek kanal bozulmalarının (çok
  yollu sönümleme, Doppler kayması, donanım empedans uyumsuzluğu) SIC
  performansı üzerindeki etkisi incelenecektir.
\item
  \textbf{PlutoSDR Entegrasyonu:} Daha erişilebilir ve düşük maliyetli
  bir SDR platformu olan Analog Devices ADALM-PlutoSDR ile test
  edilmesi.
\item
  \textbf{Çok Kullanıcılı Genişletme:} Mevcut 2 kullanıcılı yapının 3
  veya daha fazla kullanıcıya genişletilmesi ve kaskat SIC'nin hata
  yayılımı üzerindeki etkisinin araştırılması.
\item
  \textbf{Uyarlanır Güç Tahsisi:} Anlık kanal durum bilgisine (Channel
  State Information --- CSI) dayalı dinamik güç tahsis algoritmalarının
  entegrasyonu. Kanal kalitesine bağlı olarak \(a_1\) ve \(a_2\)
  katsayılarının uyarlanır biçimde güncellenmesi.
\item
  \textbf{Gelişmiş SIC Yöntemleri:} Sert karar tabanlı SIC yerine, LLR
  (Log-Likelihood Ratio) bilgisini SIC yeniden oluşturma zincirinde
  kullanan yumuşak karar SIC implementasyonu.
\item
  \textbf{OFDM-NOMA Entegrasyonu:} BPSK'dan OFDM-NOMA mimarisine geçiş
  ile frekans-seçici kanallarda çalışma kabiliyetinin kazandırılması.
\item
  \textbf{Alternatif LDPC Konfigürasyonları:} Projede mevcut olan
  \texttt{n\_0300\_k\_0152\_gap\_03.alist} dosyasındaki kısa blok LDPC
  kodunun (\(n = 300\), \(k = 152\), \(R \approx 0.507\)) düşük gecikme
  gereksinimlerindeki performansının değerlendirilmesi.
\item
  \textbf{BER Eğrisi Analizi:} Farklı SNR değerlerinde sistematik BER
  eğrilerinin çıkarılması, teorik sınırlarla karşılaştırılması ve
  SIC'nin kodlama kazancı üzerindeki etkisinin nicel olarak
  belirlenmesi.
\end{enumerate}

\begin{center}\rule{0.5\linewidth}{0.5pt}\end{center}

\subsection{7. KAYNAKÇA}\label{kaynakuxe7a}

{[}1{]} ITU-R, ``IMT Vision -- Framework and overall objectives of the
future development of IMT for 2020 and beyond,'' Recommendation ITU-R
M.2083-0, 2015.

{[}2{]} Z. Ding, Y. Liu, J. Choi, Q. Sun, M. Elkashlan, C. L. I, and H.
V. Poor, ``Application of Non-Orthogonal Multiple Access in LTE and 5G
Networks,'' \emph{IEEE Communications Magazine}, vol.~55, no. 2,
pp.~185--191, Feb.~2017.

{[}3{]} T. M. Cover, ``Broadcast Channels,'' \emph{IEEE Transactions on
Information Theory}, vol.~18, no. 1, pp.~2--14, Jan.~1972.

{[}4{]} Y. Saito, Y. Kishiyama, A. Benjebbour, T. Nakamura, A. Li, and
K. Higuchi, ``Non-Orthogonal Multiple Access (NOMA) for Cellular Future
Radio Access,'' in \emph{Proc. IEEE 77th Vehicular Technology Conference
(VTC Spring)}, Jun.~2013.

{[}5{]} L. Dai, B. Wang, Y. Yuan, S. Han, C. L. I, and Z. Wang,
``Non-Orthogonal Multiple Access for 5G: Solutions, Challenges,
Opportunities, and Future Research Trends,'' \emph{IEEE Communications
Magazine}, vol.~53, no. 9, pp.~74--81, Sep.~2015.

{[}6{]} Z. Ding, Z. Yang, P. Fan, and H. V. Poor, ``On the Performance
of Non-Orthogonal Multiple Access in 5G Systems with Randomly Deployed
Users,'' \emph{IEEE Signal Processing Letters}, vol.~21, no. 12,
pp.~1501--1505, Dec.~2014.

{[}7{]} S. M. R. Islam, N. Avazov, O. A. Dobre, and K. S. Kwak,
``Power-Domain Non-Orthogonal Multiple Access (NOMA) in 5G Systems:
Potentials and Challenges,'' \emph{IEEE Communications Surveys \&
Tutorials}, vol.~19, no. 2, pp.~721--742, 2017.

{[}8{]} R. G. Gallager, ``Low-Density Parity-Check Codes,'' \emph{IRE
Transactions on Information Theory}, vol.~8, no. 1, pp.~21--28,
Jan.~1962.

{[}9{]} D. J. C. MacKay and R. M. Neal, ``Near Shannon limit performance
of low density parity check codes,'' \emph{Electronics Letters},
vol.~32, no. 18, pp.~1645--1646, 1996.

{[}10{]} GNU Radio Project, ``GNU Radio Manual and C++ API Reference,''
{[}Çevrimiçi{]}. Erişim: https://wiki.gnuradio.org/

{[}11{]} J. G. Proakis and M. Salehi, \emph{Digital Communications}, 5th
ed., McGraw-Hill, 2008.

{[}12{]} A. Goldsmith, \emph{Wireless Communications}, Cambridge
University Press, 2005.

{[}13{]} IEEE Std 802.11n-2009, ``IEEE Standard for Information
Technology --- Telecommunications and Information Exchange Between
Systems --- Local and Metropolitan Area Networks,'' IEEE, Oct.~2009.

\begin{center}\rule{0.5\linewidth}{0.5pt}\end{center}

\subsection{8. EKLER}\label{ekler}

\subsubsection{EK-A: Sistem Değişkenleri ve Blok
Envanteri}\label{ek-a-sistem-deux11fiux15fkenleri-ve-blok-envanteri}

\paragraph{A.1. Tam Değişken
Listesi}\label{a.1.-tam-deux11fiux15fken-listesi}

% {\def\LTcaptype{none} % do not increment counter
\begin{tabular{@{}lll@{}}
\toprule
Değişken & Değer & Açıklama \\
\midrule
\bottomrule
\texttt{samp\_rate} & 200,000 Hz & Örnekleme frekansı \\
\texttt{sps} & 4 & Sembol başına örnek \\
\texttt{payload\_size} & 77 bayt & Paket yük boyutu \\
\texttt{preamble\_size} & 250 bayt & Preambül boyutu \\
\texttt{postamble\_size} & 8 bayt & Postambül boyutu \\
\texttt{noise} & 0.1 & AWGN gürültü gerilimi (ayarlanabilir) \\
\texttt{freq\_offset} & 0.01 & Normalize frekans kayması
(ayarlanabilir) \\
\texttt{time\_offset} & 1.0001 & Zamanlama sapma oranı
(ayarlanabilir) \\
\texttt{constel} & BPSK konstelasyonu & \([-1-j, -1+j, 1+j, 1-j]\) \\
\texttt{hdr} & \texttt{header\_format\_default(...)} & Başlık format
nesnesi \\
\texttt{preamble\_syms} & 64 BPSK sembol &
\texttt{{[}0xC0,\ 0xAF{]}\ ×\ 4} → bit → BPSK \\
\texttt{ldpc\_enc} & IEEE 802.11n LDPC encoder & n=1296, k=648 \\
\texttt{ldpc\_dec} & IEEE 802.11n LDPC decoder & max\_iter=50 \\
\texttt{ldpc\_dec\_2} & IEEE 802.11n LDPC decoder & max\_iter=50 (User
2) \\
\end{tabular
}

\paragraph{A.2. Blok Envanteri
Özeti}\label{a.2.-blok-envanteri-uxf6zeti}

% {\def\LTcaptype{none} % do not increment counter
\begin{tabular{@{}
  >{\raggedright\arraybackslash}p{(\linewidth - 4\tabcolsep) * \real{0.3333}}
  >{\raggedright\arraybackslash}p{(\linewidth - 4\tabcolsep) * \real{0.3333}}
  >{\raggedright\arraybackslash}p{(\linewidth - 4\tabcolsep) * \real{0.3333}}@{}}
\toprule
\begin{minipage}[b]{\linewidth}\raggedright
Kategori
\end{minipage} & \begin{minipage}[b]{\linewidth}\raggedright
Sayı
\end{minipage} & \begin{minipage}[b]{\linewidth}\raggedright
Örnekler
\end{minipage} \\
\midrule
\bottomrule
Kaynak/Çıkış Blokları & 8 & 2 File Source, 2 File Sink, 4 Vector
Source \\
Değişken Blokları & 14 & samp\_rate, sps, payload\_size, \ldots{} \\
FEC Blokları & 5 & 3 Encoder, 2 Decoder \\
İşleme Blokları & \textasciitilde30 & Repack, Scrambler, CRC, Mux,
\ldots{} \\
Özel Python Blokları & 3 & 2 Soft Diff Decoder, 1 SIC Aligner \\
Sanal Yönlendirme & 32 & 16 Sink + 16 Source \\
GUI Blokları & 6 & 3 Time Sink, 3 Range Slider \\
\textbf{Toplam} & \textbf{\textasciitilde60+} & --- \\
\end{tabular
}

\subsubsection{EK-B: LDPC Kod Parametreleri
Karşılaştırması}\label{ek-b-ldpc-kod-parametreleri-karux15fux131laux15ftux131rmasux131}

Projede iki adet LDPC parite kontrol matrisi bulunmaktadır:

% {\def\LTcaptype{none} % do not increment counter
\begin{tabular{@{}
  >{\raggedright\arraybackslash}p{(\linewidth - 4\tabcolsep) * \real{0.3333}}
  >{\raggedright\arraybackslash}p{(\linewidth - 4\tabcolsep) * \real{0.3333}}
  >{\raggedright\arraybackslash}p{(\linewidth - 4\tabcolsep) * \real{0.3333}}@{}}
\toprule
\begin{minipage}[b]{\linewidth}\raggedright
Parametre
\end{minipage} & \begin{minipage}[b]{\linewidth}\raggedright
\texttt{n\_1296\_k\_0648\_ieee.alist} (Aktif)
\end{minipage} & \begin{minipage}[b]{\linewidth}\raggedright
\texttt{n\_0300\_k\_0152\_gap\_03.alist} (Yedek)
\end{minipage} \\
\midrule
\bottomrule
Kod sözcüğü uzunluğu (\(n\)) & 1296 & 300 \\
Mesaj uzunluğu (\(k\)) & 648 & 152 \\
Kod oranı (\(R\)) & \(0.500\) & \(\approx 0.507\) \\
Standart & IEEE 802.11n & Özel tasarım \\
Alt-matris boyutu (\(Z\)) & 54 & --- \\
Yapı & Yarı-döngüsel (quasi-cyclic) & Düzensiz \\
Dosya boyutu & 44,030 bayt & 7,950 bayt \\
Kullanım & Ana sistem (NOMA + LDPC) & Alternatif / test \\
\end{tabular
}

\subsubsection{EK-C: Dosya Listesi ve
Açıklamaları}\label{ek-c-dosya-listesi-ve-auxe7ux131klamalarux131}

% {\def\LTcaptype{none} % do not increment counter
\begin{tabular{@{}
  >{\raggedright\arraybackslash}p{(\linewidth - 4\tabcolsep) * \real{0.3333}}
  >{\raggedright\arraybackslash}p{(\linewidth - 4\tabcolsep) * \real{0.3333}}
  >{\raggedright\arraybackslash}p{(\linewidth - 4\tabcolsep) * \real{0.3333}}@{}}
\toprule
\begin{minipage}[b]{\linewidth}\raggedright
Dosya
\end{minipage} & \begin{minipage}[b]{\linewidth}\raggedright
Boyut
\end{minipage} & \begin{minipage}[b]{\linewidth}\raggedright
Açıklama
\end{minipage} \\
\midrule
\bottomrule
\texttt{NOMA.grc} & 77,375 B & Ana NOMA-SIC akış diyagramı (GRC) \\
\texttt{NOMA.py} & 34,987 B & GRC'den üretilen Python betiği (627
satır) \\
\texttt{NOMA\_epy\_block\_0.py} & 3,191 B & Yumuşak diferansiyel çözücü
--- User 1 (104 satır) \\
\texttt{NOMA\_epy\_block\_0\_0.py} & 3,191 B & Yumuşak diferansiyel
çözücü --- User 2 (104 satır) \\
\texttt{NOMA\_epy\_block\_1.py} & 7,817 B & SIC hizalama bloğu ---
basic\_block (170 satır) \\
\texttt{LDPC.grc} & 36,621 B & Tek kullanıcılı LDPC test akış
diyagramı \\
\texttt{LDPC.py} & 22,904 B & LDPC test betiği (477 satır) \\
\texttt{LDPC\_epy\_block\_0.py} & 3,191 B & Yumuşak diferansiyel çözücü
--- LDPC test \\
\texttt{n\_1296\_k\_0648\_ieee.alist} & 44,030 B & IEEE 802.11n LDPC
matrisi (aktif) \\
\texttt{n\_0300\_k\_0152\_gap\_03.alist} & 7,950 B & Kısa blok LDPC
matrisi (yedek) \\
\texttt{Bpsk\_file\_transfer5.grc} & 26,215 B & Temel BPSK test
(LDPC'siz) \\
\texttt{Bpsk\_file\_transfer\_loopback.grc} & 37,983 B & USRP E310
geri-döngü testi (2.435 GHz) \\
\texttt{noma\_sic\_aligner.py} & 9,074 B & Alternatif SIC hizalama
modülü (etiket tabanlı) \\
\texttt{capture\_sic\_amplitude.py} & 740 B & SIC genlik tanılama
betiği \\
\texttt{bpsk\_transmit.txt} & 2,310 B & User 1 verici kaynak dosyası \\
\texttt{bpsk\_transmit\_2.txt} & 2,310 B & User 2 verici kaynak
dosyası \\
\texttt{bpsk\_receive.txt} & 2,310 B & User 1 alıcı çıkış dosyası \\
\texttt{bpsk\_receive\_2.txt} & 1,540 B & User 2 alıcı çıkış dosyası \\
\texttt{debug\_sic.txt} & 3,825 B & SIC hata ayıklama günlüğü \\
\end{tabular
}

\subsubsection{EK-D: Sinyal Akış Diyagramı Bağlantı
Haritası}\label{ek-d-sinyal-akux131ux15f-diyagramux131-baux11flantux131-haritasux131}

\paragraph{D.1. Verici Tarafı:}\label{d.1.-verici-tarafux131}

\begin{verbatim}
+--------------------------- User 1 TX Chain --------------------------------+
| File Source(tx1) -> Stream2TS(77B) -> CRC32 -> Repack(8->1) -> Scrambler ->    |
| LDPC Enc -> Multiply Length(×2) -> [Header Formatter + LDPC Payload] ->      |
| TagMux[Preamble(250B), Header, LDPC_Payload, Postamble(8B)] ->             |
| Constellation Mod(DBPSK, sps=4, b=0.35) -> Tag Gate ->                     |
| Multiply Const(a1 = 0.894) --------------------------------------+        |
+------------------------------------------------------------------------------+
                                                                    |
                                                                    v
                                                        +--- Add (complex) ---+
                                                        |  Süperpozisyon     |
                                                        |  s = a1x1 + a2x2  |---> Channel Model
                                                        +---------------------+
                                                                    ^
+--------------------------- User 2 TX Chain --------------------------------+
| File Source(tx2) -> Stream2TS(77B) -> CRC32 -> Repack(8->1) -> Scrambler ->    |
| LDPC Enc -> ... (aynı yapı) ... -> Tag Gate ->                               |
| Multiply Const(a2 = 0.447) --------------------------------------+        |
+------------------------------------------------------------------------------+
\end{verbatim}

\paragraph{D.2. Alıcı Tarafı
(SIC):}\label{d.2.-alux131cux131-tarafux131-sic}

\begin{verbatim}
Channel ---> FFT RRC ---> SymSync ---> CostasLoop ---> Corr Estimator
Model       (b=0.35)    (sps=4->1)   (ω=2π/100)     (preamble_syms)
                                         |                  |
                           +-------------+                  |
                           | "recovery"                     | "recovery_2"
                           v                                v
                    Soft Const Dec               +---- SIC Aligner (epy_block_1) ----+
                           |                     |  Port 0: recovery_2 (alınan)       |
                    Soft Diff Dec                |  Port 1: user_2_sub (yeniden oluş.)|
                    (epy_block_0)                |  Çıkış: temizlenmiş User 2 sinyali |
                           |                     +------------------------------------+
                    Correlate AC                            |
                           |                                v
                    LDPC Decoder ◄-----------     Soft Const Dec (2.)
                    (ldpc_dec)                            |
                           |                       Soft Diff Dec
                    Length ×0.5                     (epy_block_0_0)
                           |                              |
                    Descrambler                    Correlate AC (2.)
                           |                              |
                    Repack (1->8)                   LDPC Decoder (ldpc_dec_2)
                           |                              |
                    CRC32 Check                    Length ×0.5 -> Descramble ->
                      |     |                      Repack -> CRC32 Check (2.)
               File Sink  SIC Re-encode                    |
               (rx1.txt)    |                         File Sink
                            v                         (rx2.txt)
                    LDPC Re-encode (enc_ldpc_sic)
                            |
                    Chunks to Symbols (BPSK: [-1, +1])
                            |
                    Multiply Const (0.894)
                            |
                            v
                    "user_2_sub" -> SIC Aligner Port 1
\end{verbatim}

\begin{center}\rule{0.5\linewidth}{0.5pt}\end{center}

\emph{Bu rapor, GNU Radio Companion ortamında tasarlanan NOMA-SIC dosya
transferi sisteminin eksiksiz teknik dokümantasyonunu sunmaktadır. Tüm
blok parametreleri, matematiksel modeller ve sinyal akış yolları, kaynak
kodlar (\texttt{NOMA.grc}, \texttt{NOMA.py},
\texttt{NOMA\_epy\_block\_*.py}) incelenerek derlenmiştir.}

\end{document}
